% =====================================================================
%  Anonymity as an Aggregation Condition
%  Working paper. LaTeX is the single source of truth for THIS DOCUMENT.
%  The book (manuscript/*.md) is the single source of truth for the
%  substantive claims; every figure below is traceable to
%  model/book-calculations.ipynb and appendix/APPENDIX.md.
%
%  Build:  pdflatex anonymity-as-an-aggregation-condition.tex   (x2)
% =====================================================================
\documentclass[11pt,a4paper]{article}

\usepackage[T1]{fontenc}
\usepackage[utf8]{inputenc}
\usepackage{lmodern}
\usepackage{microtype}
\usepackage[margin=1in]{geometry}
\usepackage{amsmath,amssymb}
\usepackage{booktabs}
\usepackage{longtable}
\usepackage{array}
\usepackage{caption}
\usepackage{enumitem}
\usepackage{titlesec}
\usepackage[hidelinks]{hyperref}

\captionsetup[table]{font=small,labelfont=bf,skip=4pt}
\setlist{itemsep=2pt,parsep=0pt,topsep=4pt}
\titleformat{\section}{\normalfont\large\bfseries}{\thesection.}{0.6em}{}
\titleformat{\subsection}{\normalfont\normalsize\bfseries}{\thesubsection}{0.6em}{}

\newtheorem{proposition}{Proposition}
\newtheorem{corollary}[proposition]{Corollary}
\newtheorem{remark}{Remark}

\setlength{\emergencystretch}{2em}

\newcommand{\smallnote}[1]{\par\smallskip{\small #1}\par}

\title{\bfseries Anonymity as an Aggregation Condition:\\[0.3em]
\large Governance, Resource Structure, and Membership Dynamics\\
in Twelve-Step Mutual-Aid Organizations}
\author{\normalsize\textit{Working paper. Author anonymous by design; see the scope note.}}
\date{\normalsize August 2026}

\begin{document}
\maketitle

\begin{abstract}
\noindent
Alcoholics Anonymous (AA) has operated for more than ninety years as a worldwide
fellowship of autonomous groups under a written constitution, the Twelve
Traditions, that has not been substantially revised. The current project has not
read the AA service document needed to verify a present group count, so no count
is asserted here. This paper develops three
connected formal results and reports the robustness analysis that determines how
much each of them is worth.

First, and centrally, three of the Traditions (group conscience with servant
leadership, the refusal to organize hierarchically, and anonymity with principles
placed before personalities) are read as jointly enforcing the condition that
maximum influence weight within a group vanishes as the group grows, which Golub
and Jackson (2010) prove is necessary and sufficient for naive opinion averaging
to aggregate dispersed private information. If the reading is right, AA's
governance rules implement a known aggregation criterion, reached inductively in
1946, six decades before its formal statement. \emph{The reading is the paper's
largest unverified step and is flagged as such wherever it is used.} Two
corollaries not stated in the Traditions follow: rotation of service satisfies the
criterion only if the rotation pool scales with group size, and whether a
concentration of attention obstructs learning at all depends on how it scales with
the group rather than on how severe it looks at any one size.

Second, we derive the coupling between the Twelve Steps and Twelve Traditions
through an intermediate layer of group-produced resources, rejecting the natural
conjecture that the two lists pair by index and recovering, without assuming them,
a two-tier functional division among the Traditions and the primacy of unity that
Tradition~1 asserts of itself. We then report, against our own interest, that
five of the twelve rejections are structurally forced by sparsity and invisible to
multiplicative perturbation, that a threshold design able to see them finds them
failing by comparable margins, and that the central coupling claim survives
structural randomization in only 40.6 per cent of draws, which is less often than
it fails. The coupling results rest on the magnitudes of two hand-written
matrices and are argued for rather than certified.

Third, a stochastic simulation with endogenous membership, respecting the
institutional fact that an AA group cannot refuse \emph{membership}, distinguishes
two channels of newcomer supply and shows that groups starved of exogenous
referrals decline demographically while per-member quality among survivors is
maintained or rises, whereas groups with exclusionary internal culture shrink
without dissolving. Sensitivity analysis over all 118 registered numeric values
uses 1,002 global draws, 2,000 tiered and randomized-matrix draws, 944 multi-level
one-at-a-time points, twenty Morris trajectories, a 1,024-row Sobol base design,
and five architectural variants. In the 400-seed structural comparisons, referral
loss produces lower existence, endpoint viability and mean final membership than
pure attraction loss in all five architectures. The parameter screens report
strict, tied and reversed draws separately and do not turn that finite-domain
result into a universal claim. No parameter is estimated. The contribution is a
correspondence, a derivation, an accounting of what the analysis cannot support,
and a set of falsifiable predictions, several testable with records AA service
structures already keep.
\end{abstract}

\begin{quote}\small
\textbf{Scope and ethics.} This paper does not reproduce the text of the Twelve
Steps or Twelve Traditions, which is copyrighted by Alcoholics Anonymous World
Services, Inc.; short paraphrases are used throughout. AA is not affiliated with
this work and, by its Sixth Tradition, could not be. Nothing in this paper can
assess any individual's recovery, and it must not be used to do so. No AA
publication was acquired for this project; the historical material is drawn from
independent scholarship and from primary temperance sources in the public domain.
\end{quote}

\tableofcontents
\newpage

\section{Introduction}

Two literatures have circled Alcoholics Anonymous without meeting. The clinical
literature asks whether AA works and through what mechanism; a Cochrane review
concludes that twelve-step facilitation performs at least as well as comparison
treatments for abstinence and operates chiefly through changes in participants'
social networks (Kelly, Humphreys, and Ferri 2020). The formal-modeling literature
has produced compartmental models of drinking dynamics (S\'anchez et al.\ 2007;
Sharma and Samanta 2015), agent-based models of alcohol availability (Gorman et
al.\ 2006), and individual-level dynamical models of behavior change fitted to
clinical data (Banks et al.\ 2014, 2017). Neither literature has modeled the
content of the Twelve Steps as a structured process, and neither has modeled the
Twelve Traditions at all.

The omission is notable because the Traditions constitute an unusually clean
object of institutional study. They are a written constitution for a radically
decentralized organization: codified in 1946 from a decade of documented group
failures, formally adopted in 1950, and unamended in substance since. They govern
through autonomous groups with no enforcement mechanism, no
hierarchy, and no budget above the group level. Organizations that hold a
governance rule fixed for three-quarters of a century while scaling through six
orders of magnitude are rare, and rarer still is one whose founding documents
record the failure modes the rules were written to prevent.

\textbf{The central claim is narrow, formal, and interpretive at one joint.}
Three of the Traditions (group conscience with servant leadership, Tradition~2;
the refusal to organize hierarchically, Tradition~9; and anonymity with principles
placed before personalities, Tradition~12) share an effect that is invisible when
they are read as ethics and immediate when they are read as a constraint on a
stochastic matrix: they prevent any member from acquiring a non-vanishing share of
the group's aggregate attention. That is precisely the condition under which
DeGroot (1974) belief averaging is asymptotically wise (Golub and Jackson 2010).
To our knowledge, this correspondence has not previously been stated.

The joint is worth naming before anything is built on it. The theorem is proved
and was read at source. The mapping from three sentences of a fellowship's
constitution onto the theorem's hypothesis is a reading of those sentences,
arrived at by the author, and \emph{no computation anywhere in this paper touches
it}. Everything downstream inherits that status. Section~8.1 states what would
settle it and Section~9 states what evidence would overturn it.

Three further contributions support the central one. We derive the coupling
between Steps and Traditions from first principles, testing and rejecting the
natural conjecture that the two lists pair by index, a reading the historical
record already casts doubt on. We reformulate the Steps as a multistage production
technology in the form of Cunha, Heckman, and Schennach (2010), under which the
informal rule that steps cannot be skipped becomes the Leontief limit of a CES
family and its strictness a single estimable parameter. And we build a stochastic
membership model disciplined by an institutional constraint central to AA's
design: under Tradition~3, an AA group cannot refuse \emph{membership}, so
``closing the door'' in the sense of excluding a person from the fellowship is not
an action available to any group, and open-door adherence operates through
newcomer retention rather than newcomer arrival. The constraint proves
informative, separating empirically distinguishable channels of group decline.

\subsection{Roadmap}

Section~2 reviews the relevant literatures. Section~3 states the aggregation
framework and the mapping from Traditions to its conditions, with the
rotation-scaling and obstruction-scaling corollaries and the touring-speaker
structure. Section~4 derives the Step--Tradition coupling and reports six
robustness designs against it. Section~5 presents the CES formulation of the Steps
and its identification strategy. Section~6 describes the simulation model, its
results, and the sensitivity analysis that bounds them. Section~7 gives the
comparative historical case. Section~8 collects limitations; Section~9 states
predictions, marking where the analysis above already bears on them; Section~10
concludes. Appendix~A restates the paper in plain language. Appendix~B is a
reproducibility note.

\section{Literature Review}

\subsection{Effectiveness and mechanisms of AA}

The Cochrane review of twelve-step facilitation (Kelly, Humphreys, and Ferri 2020)
finds manualized TSF at least as effective as comparison treatments for continuous
abstinence and, in several trials, superior. The mechanism literature converges on
social network change: Kaskutas, Bond, and Humphreys (2002) find that AA's effect
on drinking outcomes is mediated substantially by changes in the composition of a
participant's social network; Rynes and Tonigan (2012) examine whether sponsorship
effects reduce to network effects. Measurement instruments for affiliation and
involvement are established (Humphreys, Kaskutas, and Weisner 1998; Tonigan,
Connors, and Miller 1996; Greenfield and Tonigan 2013). The helper-therapy
principle (Riessman 1965) and its AA-specific test (Pagano et al.\ 2004) address
whether helping others benefits the helper. Galanter (1981) supplies an earlier
account of why large-group affiliation relieves distress, which is the closest
antecedent to this paper's treatment of the group as a producer of resources
members consume.

\subsection{Formal models of drinking and recovery}

Compartmental epidemic-style models treat drinking as transmissible (S\'anchez et
al.\ 2007; Sharma and Samanta 2015). Agent-based work has modeled drinking in
relation to alcohol availability and outlet density (Gorman et al.\ 2006).
Individual-level dynamical models fitted to clinical data are due to Banks and
coauthors (2014, 2017). The relapse literature supplies the nonlinear,
multiple-equilibrium structure this paper's individual model borrows (Hufford et
al.\ 2003; Witkiewitz and Marlatt 2004, 2007).

\subsection{Skill formation}

The Steps are treated here as a multistage technology in the sense of Cunha and
Heckman (2007) and Cunha, Heckman, and Schennach (2010), with self-productivity,
dynamic complementarity, and a CES aggregator whose substitution parameter carries
the substantive content. Ben-Porath (1967) is the antecedent. Identification of
latent skill from noisy proxies follows Schennach (2004) and Hu and Schennach
(2008).

\subsection{Social learning and the wisdom of groups}

DeGroot (1974) gives the averaging model. Golub and Jackson (2010) give the
condition under which such averaging is asymptotically wise, together with the
three obstructions to it: prominent agents receiving non-vanishing attention,
imbalance between attention given and received, and insufficient dispersion
between subgroups. This paper's central section is an application of that result
and nothing more; the theorem is not extended.

\subsection{Economics of religion, clubs, and teams}

Iannaccone (1992) explains costly requirements as screening devices against
free-riding in collectives; Lembke (n.d.) applies the frame to AA directly.
Holmstr\"om (1982) gives the team production problem that a group producing a
non-excludable good faces. Ostrom (1990) supplies the design principles for
self-governing common-pool institutions, of which the Traditions are an unusually
pure instance. Angrist (2014) and Carrell, Sacerdote, and West (2013) supply the
warning this paper takes most seriously about interventions built on measured peer
effects; see Section~6.7.

\subsection{Nineteenth-century mutual-aid temperance}

The comparative case in Section~7 rests on primary temperance sources read at
source rather than on the secondary literature that descends from AA's own
account: Grosh (1842), the Washingtonian movement's own pocket manual; the
autobiographies and histories of Hawkins (1862), Marsh (1866), Gough (1869), Eddy
(1887), and Blair (1888); and the later scholarly treatments of Fehlandt (1904),
Crothers (1911), and Krout (1925). Maxwell (1950) is the one sociological
comparison of the two fellowships we have found. Kurtz (1991) is the standard
scholarly history of AA and is used for the Traditions' drafting history.

\subsection{The gap}

No prior work, so far as we have found, treats the Traditions as a formal
constraint on an influence structure, derives the Step--Tradition coupling rather
than asserting it, or models group-level membership dynamics under the
institutional constraint that a group cannot refuse membership. The gap is the
paper's occasion; whether the paper fills it well is the subject of Section~8.

\section{Governance as an Aggregation Mechanism}

\subsection{Framework}

A group of $N$ members holds beliefs about a matter of collective business. Let
$A$ be the row-stochastic trust matrix, $A_{ij}$ the weight member $i$ places on
member $j$, and let beliefs update by $b(t+1) = A\,b(t)$. Under strong
connectivity and aperiodicity, beliefs converge to a consensus equal to
$s^{\top}b(0)$, where $s$, the influence vector, is the normalized left stationary
vector of $A$. Rows describe whose beliefs a member uses; columns describe direct
attention received. Neither row equality nor column sums alone are long-run
influence. Under the regularity conditions in Golub and Jackson (2010), aggregation
requires the largest stationary weight to vanish: $\max_j s_j \to 0$.

The error is available in closed form, which is how Table~\ref{tab:error} is
computed rather than simulated. Assume here that the initial errors are iid
Gaussian with standard deviation $\sigma$. Then the consensus is
$\mu + \sum_j s_j e_j$, a normal
variable with mean zero and standard deviation $\sigma\lVert s\rVert$, and for a
mean-zero normal the expected absolute value is its standard deviation times
$\sqrt{2/\pi}$. So
\begin{equation}
\mathbb{E}\,\lvert \text{consensus} - \mu \rvert \;=\; \sigma\,\lVert s\rVert\,\sqrt{2/\pi}.
\label{eq:err}
\end{equation}
Under equal weighting $s_j = 1/N$, so $\lVert s\rVert = N^{-1/2}$ and the error is
exactly $\sigma\sqrt{2/\pi}\,/\sqrt{N}$. The single-member baseline is the same
expression at $N=1$: $\sigma\sqrt{2/\pi} = 0.798$ at $\sigma = 1$. Every figure in
this section is deterministic algebra under that Gaussian benchmark, not Monte
Carlo, and carries no sampling error. Independence and a common variance without
the Gaussian premise are not sufficient for Equation~\ref{eq:err}.

\subsection{The mapping}

\begin{table}[htbp]\centering\small
\begin{tabular}{@{}p{0.30\textwidth}p{0.32\textwidth}p{0.30\textwidth}@{}}
\toprule
\textbf{Tradition (paraphrase)} & \textbf{Formal content} & \textbf{Role} \\
\midrule
2. Group conscience; leaders serve, do not govern & Offices are hypothesized not to confer persistent attention received & Proposed mechanism reducing stationary concentration \\
9. No hierarchy; service rotates & Office-linked attention is periodically reassigned across a broad pool & Proposed mechanism preventing persistent stationary concentration \\
12. Anonymity; principles before personalities & Some status cues on which attention may condition are suppressed & Proposed mechanism reducing one source of concentration \\
\addlinespace
1. Common welfare first & $A$ strongly connected & Precondition for convergence \\
4. Group autonomy & No cross-group influence aggregation & No prominent group at higher levels \\
3. Desire to stop drinking is the only requirement for membership & $N$ unbounded; no screening on membership & Makes the asymptotic regime the relevant one \\
\bottomrule
\end{tabular}
\caption{The mapping. \textbf{This table is an interpretation of the Traditions'
wording, not a result.} It is the paper's central and least verified step.}
\label{tab:mapping}
\end{table}

\begin{proposition}[Vanishing stationary influence]
\label{prop:one}
For a sequence of row-stochastic influence matrices satisfying the convergence
and signal conditions in Golub and Jackson (2010), beliefs aggregate if the
largest normalized stationary influence vanishes, $\max_j s_j\to0$. Persistent
stationary concentration on a bounded set is sufficient for failure. A doubly
stochastic matrix is a special case: its stationary vector is uniform and
$s_j=1/N$.
\end{proposition}

The institutional claim is a separate hypothesis: Traditions 2, 9, and 12 may
make persistent stationary concentration less likely. Their wording does not by
itself imply exchangeability or double stochasticity. Conversely, violating a
Tradition is not sufficient for aggregation failure; names, offices, and hierarchy
can exist while $\max_j s_j$ still vanishes.

\begin{remark}[The status of Proposition~\ref{prop:one}]
The established mathematics and the proposed mapping must not be merged. Reading
``leaders are trusted servants; they do not govern'' as a mechanism limiting
attention received is an interpretation of a sentence, and a reader may
reasonably hold that Tradition~2 is about humility rather than weighting, or that
anonymity is chiefly protective of individuals and only incidentally structural.
Nothing computed anywhere in this paper validates that reading. Exchangeability
is an additional assumption, not a consequence of the absence of a named office.
The testable institutional hypothesis is directional: the Traditions reduce
persistent concentration in $s$, not that they make every realized weight equal.
\end{remark}

\subsection{Failure modes, quantified}

Tables~\ref{tab:influence} and \ref{tab:error} report the two quantities
separately. They are distinct and are easily conflated: at $N = 10$ the flat
regime's maximum influence weight is 0.100 and its consensus error is 0.252, and
only the first is bounded by construction.

\begin{table}[htbp]\centering\small
\begin{tabular}{@{}rrrrrr@{}}
\toprule
$N$ & Flat & Dominant & Caucus of 3 & Closed core of 5 & Rotating, pool 12 \\
\midrule
10  & 0.100 & 0.350 & 0.167 & 0.200 & n/a \\
50  & 0.020 & 0.350 & 0.167 & 0.200 & 0.041 \\
250 & 0.004 & 0.350 & 0.167 & 0.200 & 0.032 \\
500 & 0.002 & 0.350 & 0.167 & 0.200 & 0.030 \\
\bottomrule
\end{tabular}
\caption{Maximum influence weight $\max_j s_j$. Constructions: \emph{dominant},
one member receives 0.35 of every row, remainder split evenly; \emph{caucus},
three members receive 0.50 of every row between them; \emph{closed core}, five
members receive 0.45 of every row and distribute their own attention only among
themselves; \emph{rotating}, one of $R=12$ members holds share 0.35 in each term,
time-averaged over the cycle (undefined at $N=10$, where the pool exceeds the
group). Exact to the digits shown; no sampling error.}
\label{tab:influence}
\end{table}

\begin{table}[htbp]\centering\small
\begin{tabular}{@{}rrrrrr@{}}
\toprule
$N$ & Flat & Dominant & Caucus of 3 & Closed core of 5 & Rotating, pool 12 \\
\midrule
10  & 0.252 & 0.328 & 0.275 & 0.357 & n/a \\
50  & 0.113 & 0.289 & 0.238 & 0.357 & 0.132 \\
250 & 0.050 & 0.281 & 0.232 & 0.357 & 0.093 \\
500 & 0.036 & 0.280 & 0.231 & 0.357 & 0.087 \\
\bottomrule
\end{tabular}
\caption{Mean $\lvert$consensus $-$ truth$\rvert$ at $\sigma = 1$, from
equation~(\ref{eq:err}), against a single-member baseline of 0.798. Same
constructions as Table~\ref{tab:influence}.}
\label{tab:error}
\end{table}

Under flat weighting, error declines as exactly $N^{-1/2}$; under a dominant
member or an entrenched caucus it plateaus. The analytic limits confirm the
tables: dominant tends to $0.35\sqrt{2/\pi} = 0.279$; the caucus tends to 0.230;
the closed core sits at exactly $\sqrt{2/\pi}/\sqrt{5} = 0.357$, the error of a
five-member group, because the influence of every member outside the core is
identically zero at every $N$, which is why that column does not vary with $N$ at
all.

The closed core is qualitatively worse than the other two and we did not expect
that before computing it. The first two put a floor under the group's error. The
third deletes the rest of the group from the calculation: in the limit all
influence accrues to the closed subgroup and the group's accuracy is exactly that
of the subgroup alone, however many other people are present. This is a direct
consequence of Golub and Jackson's imbalance condition and standard for absorbing
sets in Markov chains; what is worth reporting is the magnitude.

Consensus is reached in every regime. Concentration does not produce visible
dysfunction, only degraded accuracy delivered with undiminished confidence. Both
of the first two pathologies have vernacular names in AA: the old-timer whose view
settles every group conscience, and the caucus that has decided before the
business meeting convenes.

\smallnote{\textbf{Severity is not the point, and it does not take much.} The
illustration uses $\alpha = 0.35$, which is a great deal, and a reader could take
from it that only substantial dominance matters. A member holding five per cent of
every row floors the consensus error at 0.0465 against a flat 0.0252 at $N=1000$,
a factor of 1.84, and that factor grows without bound as the group grows because
the flat error keeps falling and the other does not. At ten per cent the factor is
3.3.}

\subsection{Obstruction is a property of scaling, not of severity}

Golub and Jackson's three obstructions (prominent agents, imbalance, insufficient
dispersion) map naturally onto the dominant old-timer, the member whose
sponsorship lineage listens to him but who listens to nobody, and the clique.
\textbf{That mapping was asserted before it was computed, and computing it changed
it.} The three obstructions are properties of \emph{sequences} of societies as
they grow, not properties of a room.

Hold the magnitudes fixed and neither the clique nor the imbalanced member
obstructs anything. A clique of three giving a tenth of its attention outward
holds, in aggregate, 0.750 of the influence at $N=10$, 0.375 at $N=50$, 0.107 at
$N=250$ and 0.029 at $N=1000$. A member receiving twenty times the attention he
gives holds 0.690 at $N=10$, 0.290 at $N=50$, 0.074 at $N=250$ and 0.020 at
$N=1000$. Both shares vanish, both satisfy the condition, and neither is an
obstruction.

Now let the same two practices scale with the group. A clique whose inwardness
approaches one as the group expands holds 0.167 at every size tested from 10 to
1000. A member whose attention advantage grows in proportion to the group runs
from 0.182 at $N=10$ to 0.168 at $N=1000$, converging to a positive share rather
than settling on one exactly. Neither falls toward zero, and that is the whole of
what makes them obstructions.

\begin{corollary}[Obstruction scaling]
Whether a concentration of attention obstructs group learning is determined by how
it scales with $N$, not by how severe it is at any one $N$. The diagnostic
question about any concentration is therefore not how large it is but whether it
would still be there if the group doubled.
\end{corollary}

It is easy to read the three obstructions as three things a room can have at a
given moment, and on that reading a tight clique is an obstruction wherever it
appears. It is not. The general form of the point is the rotation result that
follows.

\subsection{Rotation must scale}

\begin{proposition}[Rotation breadth]
If service rotates over a pool of $R$ members, each officeholder attracting
attention share $\alpha$, time-averaged maximum influence is approximately
$\alpha/R$ plus residual flat weight. With $R$ fixed as $N$ grows, $\max_j s_j$ is
bounded below and the wisdom condition fails despite rotation.
\end{proposition}

At $N=400$ and $\alpha = 0.35$, the sweep over $R$ gives maximum influence 0.118,
0.060, 0.031, 0.016, 0.009, 0.005 and 0.003 at $R = 3, 6, 12, 25, 50, 100$ and
400, against a flat benchmark of 0.0025. The $\alpha/R$ approximation tracks the
computed value closely at every point, which is the check that the mechanism is
the one claimed. Holding the pool at twelve while the group grows, maximum
influence converges to a floor (0.041, 0.035, 0.032, 0.031, 0.030 at $N = 50, 100,
200, 400, 800$) while the flat benchmark does not (0.020, 0.010, 0.005, 0.0025,
0.0013). The ratio between them runs 2.1, 3.5, 6.4, 12.3, 23.9. That divergence is
the entire result.

The pool required to come within a factor of two of flat is 13 at $N=50$, 26 at
100, 52 at 200, 104 at 400 and 208 at 800: \textbf{twenty-six per cent of the
membership throughout}, which is the operational form of the corollary.

\smallnote{\textbf{On the size of the effect, against overstating it.} A fixed
pool of twelve at $N = 400$ should not be described as approaching the caucus
regime. At 0.031 maximum influence and 0.089 error it sits between the flat and
concentrated cases and is closer to flat on both measures, against a caucus of
three at 0.167 and 0.231. It is a real but moderate permanent cost, not a failure
of the same kind as a caucus.}

\subsection{Rotation may be aimed at the wrong positions}
\label{sec:elders}

The rotation result prices a pool that is too narrow. A separate objection, raised
by the fellowship's own commentary, is that the pool may be irrelevant.
\emph{Twelve Steps and Twelve Traditions} (AAWS 1953) describes a group's
rotating committee as sharply limited in
authority, states that in no sense whatever can its members govern or direct the
group, and then locates leadership elsewhere: in ``elder statesmen,'' former
officeholders who are called the real and permanent leadership, who are said to
become the voice of the group conscience, and to whom a perplexed group inevitably
turns. Those members hold no office and therefore rotate out of nothing.

We model this directly. Let $e$ members hold share $\alpha_e$ of every row
collectively, with no rotation, and let the remaining attention be flat. This is
not the dominant-agent family reparameterised: the share is collective and there
is no cycle to time-average.

\begin{proposition}[Advisory concentration]
Under the structure above, $\max_j s_j \to \alpha_e / e$ as $N \to \infty$, which
is bounded away from zero for any $\alpha_e > 0$. The wisdom condition fails
independently of the rotation pool $R$, since the concentration lies outside it.
\end{proposition}

With $e = 3$ and $\alpha_e = 0.10$, maximum influence is 0.1233, 0.0513, 0.0369 and
0.0342 at $N = 10, 50, 250$ and 1000, against the predicted floor of 0.0333 and a
flat benchmark falling to 0.0010. Expected consensus error is 0.2552, 0.1214,
0.0681 and 0.0525 against flat values of 0.2523, 0.1128, 0.0505 and 0.0252: a
penalty of one per cent at $N = 10$ rising to a factor of 2.08 at $N = 1000$.

The decision-relevant contrast is whether the twenty-six per cent corollary
rescues such a group. It does not. Fixing $R = \lceil 0.26 N \rceil$ at every size
and reporting $\max_j s_j$ as a multiple of the flat benchmark: rotation alone
holds at 2.0, 2.0, 2.0 at $N = 50, 250, 1000$, which is a self-consistency check
since twenty-six per cent was defined as the within-a-factor-of-two pool. Adding
three elders at $\alpha_e = 0.10$ gives 2.5, 9.2, 34.2. At $\alpha_e = 0.20$,
4.1, 17.4, 67.4. With five elders at $\alpha_e = 0.20$, 2.7, 10.7, 40.7.

\smallnote{\textbf{Scope, and what these numbers are not.} The parameters are
illustrative, not estimated. No source counts how many members a group defers to
on difficult questions or measures the attention they receive, so $e = 3$ and
$\alpha_e = 0.10$ were chosen to be conservative rather than fitted. The
qualitative result is parameter-free, since $\alpha_e/e > 0$ for any positive
$\alpha_e$; every specific multiple is not. Note also that this is a claim about a
practice AA's commentary describes and endorses, not about the Traditions, which
continue to satisfy the criterion in Proposition~\ref{prop:one}. The two are
consistent, and the distinction is the paper's recurring one between what a rule
specifies and what a fellowship does.}

\subsection{Touring speakers: a closed form}

A movement's attention structure can be modeled as bipartite: flat attention
within local societies, plus cross-society attention flowing through touring
speakers. Let each member give fraction \emph{out} of their attention to the
speakers and let the speakers return fraction \emph{back} of theirs to the general
membership. Then
\begin{equation}
\text{speakers' share of total influence} \;=\; \frac{\text{out}}{\text{out} + \text{back}},
\label{eq:speakers}
\end{equation}
exact to twelve decimal places, \textbf{independent of the number of members and
of the number of speakers}. Five speakers receiving three tenths of the movement's
attention and returning two tenths of their own hold six tenths of the influence
in a movement of ten and in a movement of a thousand alike.

This is stronger than the argument it replaces. It is not that the influence
vector converges on the speakers as the movement grows; the speakers' share is
fixed by a ratio and does not move with size at all, because growth adds members
to the denominator of the local channel and to the numerator of the speaker
channel in equal measure. Section~7 applies it.

\smallnote{\textbf{A degenerate case, and why the parameter \emph{back} is not
decoration.} With $\text{back} = 0$ the speakers form a closed communicating
class, the chain is not strongly connected, Golub and Jackson's theorem does not
apply, and the speakers hold \emph{all} the influence at every size and every
level of \emph{out}, including 0.1. That is a true statement about an absorbing
set and a useless one about a fellowship. The first version of this construction
had that defect and reporting it would have been a vacuous-robustness error.}

\section{Deriving the Step--Tradition Coupling}

\subsection{A natural conjecture, tested}

\label{sec:conjecture}
Two ordered lists of twelve, printed in a single volume, invite the conjecture
that Step $i$ draws on Tradition $i$. The conjecture merits testing rather than
adoption. The historical record gives grounds for doubt: the Traditions were
codified from 1946 essays written roughly a decade after the Steps, in response to
specific organizational crises, and no AA doctrine pairs the lists by index.

That last clause was previously an assertion about literature the project had not
read. It has since been checked. \textit{Twelve Steps and Twelve Traditions}
(AAWS 1953) is the only work treating both lists at length, one chapter each, by
the same author, and is therefore where an intended pairing would surface. Across
its twelve Tradition chapters, none cites the Step of its own number; across its
twelve Step chapters, the word \emph{Tradition} does not occur at all. The single
indexed cross-reference in the volume runs off-index: the chapter on Tradition~8
mentions the Twelfth Step, to distinguish paid service work from twelfth-step work
itself. The count is reproducible in minutes from the files AAWS publishes free,
and the procedure is stated here rather than only its result: in each Tradition
chapter, count references to the same-numbered Step; in each Step chapter, count
occurrences of \emph{Tradition}. The
stakes are concrete. Under index-pairing, Step~5 (telling one's inventory to
another person) would depend on Tradition~5 (primary purpose), though what the
step plainly requires is confidentiality, which is Tradition~12; and Step~12
(carrying the message) would couple to anonymity, though Tradition~5 states nearly
the same sentence. We therefore derive the coupling from first principles and test
the conjecture against the result.

\subsection{Method}

To avoid re-deriving the same mistake, no direct step-to-tradition mapping is
permitted. An intermediate layer of eight group-produced resources is introduced
(admission and standing, identification, living proof, confidential audience,
counsel, a recipient for twelfth-step work, continuity, and normative pressure),
and the coupling is computed as
\begin{equation}
B \;=\; S\,G^{\top},
\end{equation}
where $S[i,r]$ records how much Step $i$'s execution consumes resource $r$
(written by asking what each step requires) and $G[j,r]$ records how much
Tradition $j$ governs the supply of $r$ (written independently of $S$). Both
matrices are hand-written judgments. That is the whole of their evidentiary
status, and Section~4.5 does not soften it.

\subsection{Results}

Table~\ref{tab:coupling} gives the unperturbed coupling.

\begin{table}[htbp]\centering\small
\begin{tabular}{@{}llrlrlrr@{}}
\toprule
Step & Principal & Value & Runner-up & Value & Index-mate & Value & Its rank \\
\midrule
1 admit        & T3  & 1.22 & T1 & 0.99 & T1  & 0.99 & 2 \\
2 believe      & T11 & 0.82 & T5 & 0.72 & T2  & 0.13 & 7 \\
3 decide       & T2  & 0.31 & T1 & 0.27 & T3  & 0.01 & 7 \\
4 inventory    & T1/T2 tie & 0.35 & --- & --- & T4  & 0.00 & 8 \\
5 tell someone & T12 & 1.08 & T1 & 0.65 & T5  & 0.12 & 5 \\
6 willing      & T1  & 0.25 & T2 & 0.24 & T6  & 0.00 & 7 \\
7 ask          & T1  & 0.17 & T2 & 0.13 & T7  & 0.00 & 7 \\
8 list harms   & T2  & 0.43 & T1 & 0.29 & T8  & 0.07 & 4 \\
9 amends       & T2  & 0.90 & T1 & 0.48 & T9  & 0.00 & 8 \\
10 daily       & T1  & 0.94 & T2 & 0.43 & T10 & 0.00 & 8 \\
11 connect     & T1  & 0.47 & T2 & 0.21 & T11 & 0.11 & 4 \\
12 carry it    & T5  & 1.25 & T3 & 1.10 & T12 & 0.12 & 6 \\
\bottomrule
\end{tabular}
\caption{The author-coded semantic overlap $B = SG^{\top}$, unperturbed. It is not
the executable state-update map. Competition ranks use $1+$ the count of strictly
larger entries, so exact ties share rank. Not one Step has its index-mate in the
maximizing set. Load per Tradition, exact: T1 6.52, T2 3.89,
T5 3.88, T3 2.69, T11 2.69, T12 2.62, T8 1.11, and exactly zero for T4, T6, T7,
T9 and T10.}
\label{tab:coupling}
\end{table}

\begin{enumerate}
\item \textbf{The author coding contains a two-tier division.} Five Traditions
(autonomy, non-endorsement, self-support, non-hierarchy, and no outside issues)
have identically zero rows in $G$. Those zeros were assigned by the author; the
matrix multiplication reveals their consequences but does not independently
discover or validate the classification.

\item \textbf{Unity has the largest raw semantic load under this coding.} Under
index-pairing, unity fed one step; in the authored matrices it overlaps with
resources consumed by nearly every Step. The multiplication summarizes those
judgments; it does not supply independent evidence for them.

\item \textbf{The index-pairing conjecture is rejected on all twelve counts.}
Step~5's principal supplier is Tradition~12 and Step~12's is Tradition~5, the two
correspondences a reader would guess, reversed.

\item \textbf{Classification is not always intuitive.} Tradition~8
(non-professionalism), which a functional reading might class as protective, lands
in the enabling tier with small load, because it governs the identity of the
confidential hearer, a fellow member rather than a clinician.
\end{enumerate}

\subsection{Six robustness designs, and what each can and cannot see}

\paragraph{The trivial-count decomposition, stated first because it deflates the
headline.} T4, T6, T7, T9 and T10 have identically zero rows in $G$, so Steps 4,
6, 7, 9 and 10 have index-mate entries of exactly zero, and index-pairing cannot
hold for them under \emph{any} perturbation that preserves sparsity. Both
multiplicative designs below preserve sparsity. The consequence is checkable and
was checked: at every level, the proportion of draws in which index-pairing fails
on all twelve equals the proportion in which it fails on the seven non-trivial
Steps, to the last draw (98.75 and 98.75 at $\pm15$ per cent; 85.50 and 85.50 at
$\pm30$; 72.45 and 72.45 at $\pm50$; 65.15 and 65.15 at $\pm75$; 40.60 and 40.60
structurally). \textbf{Five of the twelve counts are not evidence, and no
multiplicative design could have made them so.}

\paragraph{Design 1, multiplicative jitter.} Every entry of $S$ and $G$ is
multiplied by an independent uniform draw on $[1-L, 1+L]$ for
$L \in \{0.15, 0.30, 0.50, 0.75\}$; 2,000 draws per level, seed 3. This represents
a reader who disagrees with the magnitudes by up to $L$ and agrees about which
cells are empty.

\paragraph{Design 2, structural randomization.} Every non-zero entry of $S$ and
$G$ is replaced by an independent uniform draw on $[0.05, 1.00]$; zeros stay zero.
2,000 draws, seed 23. This represents a reader who accepts only the pattern of
which Tradition touches which resource and rejects every magnitude we chose. A
parallel 1,000-draw screen passes randomized nonzero matrices through the full
executable model: the referral-versus-pure-attraction comparison is strict in all
1,000 draws on final membership and endpoint viability and in 999, with one tie,
on existence. Full-adherence viability in all three screening seeds holds in only
784 draws. The comparison and the absolute outcome therefore cannot borrow one
another's robustness.

\begin{table}[htbp]\centering\small
\setlength{\tabcolsep}{4pt}
\begin{tabular}{@{}lrrrr@{}}
\toprule
Design & T1 leads & Index-pairing wrong, all & Step 5 $\to$ T12 & Step 12 $\to$ T5 \\
\midrule
jitter $\pm$15\% & 100.0 [99.8, 100.0] & 98.8 [98.2, 99.2] & 100.0 [99.8, 100.0] & 89.7 [88.3, 91.0] \\
jitter $\pm$30\% & 100.0 [99.8, 100.0] & 85.5 [83.9, 87.0] & 99.5 [99.1, 99.7] & 67.3 [65.3, 69.4] \\
jitter $\pm$50\% & 98.0 [97.3, 98.5]   & 72.5 [70.5, 74.4] & 88.7 [87.2, 90.0] & 52.9 [50.7, 55.1] \\
jitter $\pm$75\% & 86.8 [85.3, 88.3]   & 65.2 [63.0, 67.2] & 71.5 [69.5, 73.4] & 43.8 [41.6, 45.9] \\
\textbf{structural} & \textbf{75.4 [73.5, 77.2]} & \textbf{40.6 [38.5, 42.8]} & \textbf{17.5 [15.9, 19.2]} & \textbf{27.6 [25.7, 29.7]} \\
\bottomrule
\end{tabular}
\caption{Percentage of draws in which each claim holds. Wilson intervals at 95 per
cent on $n = 2{,}000$. The two 100.0 entries are 2,000 of 2,000 and should be read
as ``not observed to fail'', not as certainty.}
\label{tab:coupling-robust}
\end{table}

Step~12 is the fragile row and the reason is the margin, not the level: its top
two are T5 at 1.25 and T3 at 1.10, and both Traditions govern the recipient
resource (T5 at 0.9, T3 at 0.8), so a perturbation moving them in opposite
directions flips the winner.

\paragraph{Design 3, the threshold test, which is the only one that can see the
five.} For Step $i$, let $w_i$ be the row sum of $S$ (total consumption across the
eight resources) and $b_i$ the largest entry in row $i$ of $B$ excluding the
index-mate's own entry. If Tradition $i$ governed every resource at a uniform
strength $c$, its entry in row $i$ would be $c\,w_i$, so index-pairing holds at
Step $i$ exactly when
\begin{equation}
c \;>\; c^{*}_i \;=\; b_i / w_i .
\end{equation}
This is one division per Step, carries no sampling error, and is the only design
here that can turn a structural zero into a non-zero.

\begin{table}[htbp]\centering\small
\begin{tabular}{@{}rlrrrr@{}}
\toprule
Step & Index-mate & Beats & Row sum of $S$ & $c^{*}$ & $c^{*}$ / mean live entry \\
\midrule
1 & T1 & 1.22 & 2.40 & 0.508 & 1.36 \\
2 & T2 & 0.82 & 1.70 & 0.482 & 1.29 \\
3 & T3 & 0.31 & 0.70 & 0.443 & 1.18 \\
\textbf{4} & \textbf{T4} & 0.35 & 0.80 & \textbf{0.438} & 1.17 \\
5 & T5 & 1.08 & 1.70 & 0.635 & 1.70 \\
\textbf{6} & \textbf{T6} & 0.25 & 0.60 & \textbf{0.417} & 1.11 \\
\textbf{7} & \textbf{T7} & 0.17 & 0.40 & \textbf{0.425} & 1.14 \\
8 & T8 & 0.43 & 0.80 & 0.538 & 1.44 \\
\textbf{9} & \textbf{T9} & 0.90 & 1.50 & \textbf{0.600} & 1.60 \\
\textbf{10} & \textbf{T10} & 0.94 & 1.50 & \textbf{0.627} & 1.67 \\
11 & T11 & 0.47 & 0.80 & 0.588 & 1.57 \\
12 & T12 & 1.25 & 2.40 & 0.521 & 1.39 \\
\bottomrule
\end{tabular}
\caption{The threshold test. Bold rows are the five protective Traditions, whose
index-mate entries are structurally zero and therefore invisible to
Table~\ref{tab:coupling-robust}. Exact; no sampling error.}
\label{tab:threshold}
\end{table}

The governance matrix has 35 non-zero entries of 96 cells, with mean 0.374, median
0.300, minimum 0.10 and maximum 1.00. \emph{Every} $c^{*}$ exceeds both the mean
and the median. The protective range, 0.417 to 0.627, is not contained in the
enabling range, 0.443 to 0.635; it extends below it at both ends. Mean $c^{*}$ is
0.501 for the five protective Steps, 0.531 for the seven enabling ones, and 0.518
across all twelve. The protective Steps are therefore marginally \emph{closer} to
index-pairing holding, by 0.03 on the mean threshold, which is the direction an
objector would predict and an order of magnitude smaller than the distance to the
mean live entry. Setting $c$ to the mean live entry of 0.374 for every Step
simultaneously, all twelve index-mates still lose, by margins from 0.02 at Step~7
to 0.44 at Step~5.

What the test licenses is the statement that index-pairing would require the
index-mate to govern its own Step's needs more strongly than a typical entry in
the matrix, uniformly across all twelve. What it does not license is any statement
about whether the zeros are correctly \emph{placed}: it holds the sparsity
pattern's origin fixed and prices only its consequences. It is a sensitivity
analysis of a judgment, not a test of it.

\paragraph{Design 4, the resource-list test.} Delete each of the eight resources
in turn, then merge each pair, then delete each pair: sixty-four resource lists in
all. Only Step~1 and Step~2 ever regain their index-mates, in fifteen variants and
one respectively, and every one of the fifteen Step~1 failures involves removing
or merging the admission resource, which is precisely what gives the open door its
lead there. This establishes that the list is no \emph{finer} than it needs to be.
It cannot establish that the list is fine \emph{enough}, because inventing a ninth
resource requires a judgment about what it contains and cannot be done by
rearranging the eight.

\paragraph{Design 5, the reassignment test.} Kurtz (1991) records that in some later AA
literature the concept conveyed by \emph{single-purposed} was obfuscated by
substituting \emph{unity}. If the governance matrix absorbed that semantic drift,
unity's primacy is an artifact of it. For each of the eight resources in turn, set
$G[\text{T5},r]$ to the maximum of its current value and $G[\text{T1},r]$, set
$G[\text{T1},r]$ to zero, and recompute: this transfers unity's entire governance
of one resource to singleness of purpose, which is the strongest form of the
objection, and it is exact. \textbf{Six of the eight transfers leave unity
leading}, by margins from 1.44 to 2.54. Transferring continuity gives T5 the lead
5.14 to 4.63; transferring pressure gives it 5.12 to 4.66. So unity's primacy is
conditional on two assignments and on nothing else in the matrix. Notably, the
single resource-list variant that breaks it (Design 4) drops continuity and
pressure together, which is exactly the pair this test identifies: two independent
designs agreeing on which two resources carry a result is worth more than either
alone. The passage Kurtz names as decisive is in an AA-copyright work that has not
been obtained, so the source identified as settling the question is unread.
The 1953 commentary, read in full on 10 August 2026, keeps unity and singleness
of purpose apart across its chapters on Traditions 1 and 5, which corroborates
the assignment without being the passage Kurtz cited.

\paragraph{Design 6, sparsity perturbation.} Every design above varies magnitudes
and holds the sparsity pattern fixed; a second reader would disagree about the
pattern. Flipping cells at random within the seven enabling rows (56 cells, so the
two-tier split is held fixed), 2,000 draws per row:

\begin{table}[htbp]\centering\small
\setlength{\tabcolsep}{5pt}
\begin{tabular}{@{}rrrrr@{}}
\toprule
Cells flipped & Index-pairing wrong, all & T1 leads & Step 5 $\to$ T12 & Step 12 $\to$ T5 \\
\midrule
1  & 96.3\% & 100.0\% & 98.7\% & 94.8\% \\
2  & 92.7\% & 99.7\%  & 96.3\% & 90.5\% \\
4  & 86.2\% & 96.5\%  & 93.3\% & 80.3\% \\
8  & 75.0\% & 83.7\%  & 85.4\% & 72.0\% \\
16 & 53.0\% & 50.8\%  & 70.8\% & 53.0\% \\
\bottomrule
\end{tabular}
\caption{Sparsity perturbation. Section~4 tolerates a reader differing on about
four of fifty-six enabling cells and does not tolerate one differing on sixteen.}
\label{tab:sparsity}
\end{table}

\textbf{This is a bound, not a measurement, and the distinction is the caveat.} A
random flip is not a plausible reader. Somebody who thinks self-support governs
continuity changes that cell for a reason, and their remaining cells correlate
with the reason. Random flips are harsher in respecting no reason and gentler in
not concentrating on the cells that carry the results. The sweep answers ``how
much disagreement, counted in cells'' and says nothing about which cells a reader
would choose. It also holds the two-tier split fixed, because the split is what
defines which rows are available to flip, so the more serious disagreement, a
reader who fills a protective row, cannot be represented by any perturbation
design at all. Only completed elicitation forms can answer either question, and
none has been returned; see Section~8.

\subsection{What the coupling results rest on}

Two classes of claim in this paper rest on different things and must not be
defended the same way.

The executable model's \textbf{referral-versus-pure-attraction comparison} survives
replacing every nonzero matrix magnitude with a random value while preserving
sparsity: 1,000 of 1,000 strict draws on final membership and endpoint viability,
and 999 strict plus one tie on existence. The fully adherent group's endpoint
viability is not invariant. This result concerns the executable model, not the raw
coupling statistics in this section.

The \textbf{coupling} claims of this section do not. Index-pairing fails on all
twelve in only 40.6 per cent of structurally randomized draws, which is less often
than it holds; the Step~5 inversion survives in 17.5 per cent and the Step~12
inversion in 27.6. These claims rest on the magnitudes in two matrices we wrote,
and they must be argued for rather than certified by a robustness percentage
earned elsewhere. The one exception is unity's primacy, which survives structural
randomization in 75.4 per cent of draws, the highest figure of any claim in this
section: three-quarters of the time a random matrix with the same sparsity pattern
puts unity first anyway.

A reader who thinks the two matrices are arbitrary should not be persuaded by
Section~4, and we would rather say so than borrow an executable endpoint
comparison's robustness for a semantic-coupling result that has not earned it.

\section{The Steps as a Multistage Technology}

\begin{quote}\small
\textbf{This section is a reformulation, not the simulation's update rule, and
the two must not be read as the same object.} The CES form below is a way of
\emph{stating} the ordering claim so that its strictness becomes one estimable
parameter. The dynamic model of Section~6 uses the gated-growth equation given in
\S6.2, which is not a CES aggregator. Nothing in Section~6 depends on the value of
$\rho$.
\end{quote}

Let $x_{i,t}$ denote latent practice of step $i$. Following Cunha, Heckman, and
Schennach (2010):
\begin{equation}
x_{i,t+1} \;=\; A_i \Big[\, \gamma_{i1} x_{i,t}^{\rho_i} + \gamma_{i2} x_{i-1,t}^{\rho_i}
 + \gamma_{i3} G_{i,t}^{\rho_i} + \gamma_{i4} M_t^{\rho_i} \,\Big]^{1/\rho_i},
\end{equation}
with weights summing to one, elasticity of substitution
$\sigma_i = 1/(1-\rho_i)$, $G_{i,t}$ the resource bundle from Section~4, and $M_t$
maintenance capacity built from Steps 10 to 12. Self-productivity is
$\gamma_{i1} > 0$; the cross-partial in prior-stage stock and group input is
positive throughout (from $+1.73$ at $\rho = -4$ to $+0.08$ at $\rho = 0.5$), so
group support is worth more to a member who has done the preceding work: dynamic
complementarity in Cunha and Heckman's sense.

\begin{proposition}[The ordering rule as a limit]
As $\rho_i \to -\infty$ the aggregator converges to $\min(\cdot)$, so zero
prior-stage stock forces zero output: the informal rule that steps cannot be
skipped. The rule holds for all $\rho_i \le 0$ and fails for $\rho_i > 0$, where
the group input substitutes for the missing stage (output 0.217 at $\rho = 0.3$
with prior stage at zero).
\end{proposition}

The reduction converts a widely held but untested claim of practice into a sharp
empirical question: within this CES reformulation, the sign of $\rho$ decides
whether the Steps are a chain or a menu. Estimating that sign would require an
empirical identification design; the calculation here does not estimate it.

\subsection{Identification, anchoring, and design}

Latent practice would have to be measured with a design that addresses error in
the inputs, loadings, scale, and group-input endogeneity. Schennach (2004) and Hu
and Schennach (2008) describe identification strategies under explicit
assumptions, but this paper does not implement them. Its simulation gives the four
regressors and proxy loadings to the estimator without error and adds noise only
to the outcome. Repeating that exercise shows the modest point that averaging
three rescaled output proxies improves resolution relative to one. It does not
recover a latent-variable estimator, justify three instruments per latent per
wave, or settle the sign of $\rho$.

The binding obstacle to estimation is endogeneity of the group input: groups
direct attention to struggling members and successful members attract sponsees, so
$G$ is not exogenous to member state. CHS address the analogous endogeneity of
parental investment; adapting their approach is necessary before any fit to panel
data and is not attempted here.

\section{Membership Dynamics Under the Open Door}

\subsection{The institutional constraint, restated}

Tradition~3 makes the desire to stop drinking the sole requirement for
\textbf{membership}. It is tempting to render this as ``an AA group cannot refuse
admission,'' and that overstates it. AA groups hold closed meetings as a matter of
routine, and a closed meeting restricts who is present in a room; it does not
remove anyone's membership in the fellowship. The constraint is on membership, not
on attendance, and the model is built to that constraint. The distinction was
supplied by a reader who knows the rooms, and it is the single correction this
work has received that changed the most for the fewest words.

The distinction has modeling consequences. Arrival at a meeting does not depend on
the group's welcome in the default model. Tradition~3 instead has two separately
controllable paths: its row governs four resource columns, and it adds dropout
friction weighted by inverse practice. The second is not an early-tenure variable:
the model contains no tenure or cohort state, so a long-tenured low-practice member
receives the same weight as a recent arrival at the same practice. Inflow contains
an exogenous referral floor independent of the group's attractiveness; exit also
includes a practice-independent churn term for relocation and mortality.

Two further design choices follow from the same logic. Carrying capacity is
supplied by the established core rather than the population mean, since a room of
three veterans and twenty newcomers still contains three people able to carry a
newcomer. And the maintenance-capacity gate, whose role is to model relapse in
members with something to maintain, phases in with step index rather than
throttling entry-level growth, to which it does not conceptually apply.

\subsection{Model summary}

Members occupy a twelve-dimensional practice state. Per-step growth, in full, for
member $m$ and step $i$:
\begin{equation}
\frac{dx_i}{dt} \;=\; h(m)\, a_i\, \mathrm{gate}_i\, \mathrm{peer}_i\, C^{m}_i\,\bigl(1 - x_i\bigr)
 \;-\; d_i\, x_i,
\label{eq:growth}
\end{equation}
with
\begin{align}
\mathrm{gate}_i &= x_{i-1}^{\,p}, \qquad \mathrm{gate}_1 = 1, \\
\mathrm{peer}_i &= (1-\beta_i) + \beta_i G_i, \qquad G = S_{\text{norm}} R, \\
C^{m}_i &= 1 - w(i)\,(1 - C), \qquad w(i) = 0.05 + (i-1)\tfrac{0.95}{11}, \\
d_i &= \delta_0\bigl(1 + \psi(1 - x_{i+1})\bigr) \ \ (i < 12), \qquad d_{12} = \delta_0 .
\end{align}
Maintenance capacity $C$ is identical across steps for a given member and is
defined in two parts. Own capacity is a Hill gate on $M$, the mean of that member's
Steps 10 to 12, and group support supplies a floor beneath it:
\begin{equation}
C(M) = \frac{M^{n}}{k^{n} + M^{n}}, \quad n = 3.0,\ k = 0.12;
\qquad
C = C(M) + \bigl(1 - C(M)\bigr)\,\omega\,\overline{C}, \quad \omega = 0.75,
\label{eq:capacity}
\end{equation}
where $\overline{C}$ is the mean own-capacity across living members, so a member
whose own maintenance has collapsed retains a fraction of capacity as long as the
group around them has not. Because Steps 10 to 12 sit at the heavy end of the
$w(i)$ weighting, maintenance gates its own accumulation. That architecture can be
bistable, but it is not generally so in the corrected model's own endpoint
environments: a capability-one high/low-start test separates in 7 of 400 environments
and collapses to one low state in the mean environment. Here
$w(i)$ is per-step exposure to it, rising linearly from
0.05 at Step~1 to 1.00 at Step~12, which is the formal content of the claim that
an arrival has nothing to maintain and a veteran has a great deal; and $h(m)$ is
member heterogeneity, a lognormal draw made once when a member arrives and fixed
thereafter. It is parameterized as $\exp(N(-\sigma_h^2/2,\sigma_h))$, so its
arithmetic expectation is one and changing $\sigma_h$ changes dispersion without
mechanically changing mean capability.

Group resources $R$ are produced from member states through the governance-weighted
averages of Section~4, so shortfalls do not compound multiplicatively, with
protective Traditions acting solely as multipliers on the enabling ones they
guard. Membership is endogenous: arrivals are exogenous referrals plus attraction
proportional to members' twelfth-step practice and the attraction path of
Tradition~11; exits combine practice-dependent dropout, the Tradition-3 friction
path, and churn. Tradition~11 also has a distinct resource-governance path.
Integration is Euler at $dt = 0.5$ weeks over a 1,560-week (thirty-year) horizon.
This horizon is a reporting choice, not a steady-state assertion. In 200-seed
checks, mean membership is 29.34, 21.14, 16.43 and 15.64 at 10, 20, 50 and 100
years. The 100-year condition includes one closure. Paired changes in final
membership from $dt=0.5$ to 1.0, 0.25 and 0.125 weeks are respectively 1.54
$\pm$ 1.78, 0.87 $\pm$ 1.84 and -0.41 $\pm$ 1.74. None resolves at this sample
size, so the tested step is adequate for the reported finite-horizon comparisons;
the runs do not establish an equilibrium or indefinite persistence.

Parameters: $\delta_0 = 0.06$ per week (unattended half-life 11.6 weeks),
$\psi = 0.20$, $p = 1.5$, step top speeds from 0.15 at Step~9 to 0.30 at Step~1.
The inventory is 22 continuous scalars, 12 step speeds, 49 non-zero cells in $S$
and 35 in the governance matrix: \textbf{118 registered sensitivity values out of
226 listed cells, none fitted.} The model-choice inventory separately records
fixed constants, structural zeros, equations, thresholds and experiment-design
choices, so 118 is not the count of every authored choice.

\begin{remark}[Calibration, and the exact sense in which it fails]
Inflow, dropout and churn were originally set to target a group near forty-five
members with an experienced core near nine, roughly a healthy urban meeting. That
was calibration to a stylized fact, not a dataset. After the lognormal capability
draw was corrected to have arithmetic mean one, the target fails. Across 400 runs
the model delivers 17.80 members, 95 per cent half-width 0.88. Among the 394 viable
endpoints, the \emph{established} count above 0.1 is 14.13, half-width 0.79, and the
\emph{experienced} count above 0.5 is 1.25, half-width 0.20. The two thresholds are
not interchangeable. The authored rates are not retuned after observing this
failure, and the absolute levels are not estimates of real meetings.
\end{remark}

\begin{remark}[Thirty-five parameters that cannot matter at full adherence]
The governance matrix is column-normalized, so governance quality is identically 1
when every Tradition is at 1.0 and the matrix cancels exactly. This is algebra,
not simulation. All 35 governance cells therefore produce exactly zero change in
every outcome at full adherence, and a sensitivity design that perturbs them and
reports no effect has found nothing. Thirty-five of the 118 registered values are in that
position.
\end{remark}

\subsection{Three channels of decline}

Four configurations, 400 paired seeds each, thirty-year horizon, $dt = 0.5$ weeks.
``Viable'' means more than five members; existence means at least one, and zero is
permanent closure. \textbf{Quality is established-member practice in viable groups
and is conditional throughout; the viable fraction is printed beside it in every
row.} Viability carries a 95 per cent Wilson
interval; membership and quality carry a 95 per cent half-width from the cross-run
standard error.

\begin{table}[htbp]\centering\footnotesize
\setlength{\tabcolsep}{4pt}
\begin{tabular}{@{}lrrrrrr@{}}
\toprule
Condition & Viable y10 & $N$ if viable, y10 & Quality y10 & Viable y30 & 95\% int. & Quality y30 \\
\midrule
nothing wrong & 0.9975 [0.9860, 0.9996] & $29.40 \pm 1.16$ & $0.3016 \pm 0.0052$ & 0.985 & 0.968--0.993 & $0.2648 \pm 0.0058$ \\
invisible     & 0.9925 [0.9782, 0.9974] & $12.98 \pm 0.34$ & $0.2949 \pm 0.0079$ & 0.985 & 0.968--0.993 & $0.2592 \pm 0.0071$ \\
unreferred    & 0.660 [0.612, 0.705] & $16.47 \pm 1.18$ & $0.3475 \pm 0.0084$ & 0.0275 & 0.015--0.049 & $0.3642 \pm 0.0443$ \\
unwelcoming   & 0.905 [0.872, 0.930] & $15.70 \pm 0.82$ & $0.3597 \pm 0.0097$ & 0.5475 & 0.499--0.596 & $0.3211 \pm 0.0134$ \\
\bottomrule
\end{tabular}
\caption{Decline scenarios, \textbf{400 paired seeds}, 1,560-week horizon.
\emph{Invisible} sets only Tradition~11's attraction path to zero;
\emph{unreferred} sets exogenous inflow to zero; \emph{unwelcoming} sets
both Tradition~3 paths to zero. Only eleven unreferred runs are viable at year
thirty, so its conditional quality interval is wide.}
\label{tab:decline}
\end{table}

Membership counted over all runs with closures as zero reaches, at year thirty,
$17.80 \pm 0.88$, $12.38 \pm 0.34$, $0.51 \pm 0.23$ and $6.76 \pm 0.59$ respectively.

Three signatures separate. Attraction loss leaves a smaller remnant sustained by
referrals and produces no closure. Referral loss closes 89.5 per cent of groups by
year thirty, while the eleven viable endpoints have high conditional practice.
Combined Tradition~3 loss closes 25.0 per cent and leaves 54.8 per cent viable.

The mortality profiles follow from structure rather than from hand-tuning. A group
on one engine has no floor underneath it: its inflow becomes a function of its own
state, which makes the population dynamics multiplicative rather than additive,
and a multiplicative process with no floor has an absorbing state at zero.
Invisible and unwelcoming groups retain a referral floor, but the latter can still
close through its resource and dropout paths. Unreferred groups are the only ones
whose inflow itself can go to zero and stay there.

\smallnote{\textbf{The quality claim, stated exactly.} At year ten,
established-member practice against a healthy viable group is 0.0067 lower for the
invisible condition, 0.0459 higher for the unreferred one, and 0.0581 higher for
combined Tradition~3 loss. Membership is lower in all three. Practice can therefore
mislead through selection, but the corrected model does not support the stronger
claim that the interior of a dying group offers no warning.}

\subsection{The gap between decline and death}

The unreferred condition separates two questions that a single membership series
conflates.

\begin{table}[htbp]\centering\small
\begin{tabular}{@{}lrrrr@{}}
\toprule
Series & y5 & y10 & y20 & y30 \\
\midrule
unreferred, all runs        & $21.55 \pm 1.01$ & $11.68 \pm 1.02$ & $2.50 \pm 0.48$ & $0.51 \pm 0.23$ \\
unreferred, viable only     & $22.07 \pm 1.00$ & $16.47 \pm 1.18$ & $11.39 \pm 1.37$ & $12.27 \pm 3.81$ \\
viable fraction             & 0.970 & 0.660 & 0.1725 & 0.0275 [0.015, 0.049] \\
\bottomrule
\end{tabular}
\caption{Unconditional versus viability-conditioned membership under referral loss.
Both decline sharply; conditioning still hides the mass of closures but no longer
makes membership look stable.}
\label{tab:conditional}
\end{table}

This is why conditioning is load-bearing throughout the paper. In the referral-loss
condition both the all-run and viable-only membership series decline, but the
all-run series also carries the growing mass of closures. Neither series can stand
in for the other.

\subsection{Comparative statics under variance control}

Discipline preceded comparison: baseline survival is stable across independent
seed blocks, and common random numbers reduce the paired standard error of
membership comparisons. The reference cross-seed standard deviation is 5.67,
while paired standard errors range from 0.22 to 0.33. Each Tradition is degraded
singly from 0.85 to 0.5 against a reference group of 13.10 members.

\begin{table}[htbp]\centering\small
\begin{tabular}{@{}llrrr@{}}
\toprule
Tradition & Tier & Members cost & 95\% half-width & $t$ \\
\midrule
T3 mixed adherence    & enabling   & 2.99 & 0.65 & 9.1 \\
T11 mixed adherence   & enabling   & 1.88 & 0.60 & 6.1 \\
T1 unity              & enabling   & 1.65 & 0.65 & 4.9 \\
T12 anonymity         & enabling   & 0.99 & 0.58 & 3.3 \\
T5 one purpose        & enabling   & 0.96 & 0.61 & 3.1 \\
T4 autonomy           & protective & 0.88 & 0.62 & 2.8 \\
T7 self-support       & protective & 0.88 & 0.62 & 2.8 \\
T2 group conscience   & enabling   & 0.48 & 0.62 & 1.5 \\
T6 no endorsement     & protective & 0.23 & 0.46 & 1.0 \\
T10 no outside opinion& protective & 0.23 & 0.46 & 1.0 \\
T9 no organization    & protective & 0.02 & 0.43 & 0.1 \\
T8 non-professional   & enabling   & -0.20 & 0.51 & -0.8 \\
\bottomrule
\end{tabular}
\caption{Single-Tradition degradation, \textbf{400 paired replications} under
common random numbers. Seven of twelve have intervals excluding zero. T3 and T11
are mixed adherence interventions; their paths are split in the release factorials.}
\label{tab:tradition}
\end{table}

\begin{quote}\small
\textbf{The replication budget determines the ranking, and 30 is not enough.}
Computed at 30 paired replications, this comparison returns autonomy and
self-support at 5.5 members each with $t = 2.6$ as the only pair clearing
$\lvert t \rvert > 2.5$, which invites the inference that the protective
Traditions lead the ranking, consistent with their derived role as guards on
everything else. Under the corrected model, at 400 replications mixed Tradition~3
leads at 2.99 members, mixed Tradition~11 follows at 1.88, and seven comparisons
resolve rather than two. \textbf{A common-random-numbers design at 30 replications is more
efficient than 30 independent runs and is still 30 replications.} Variance
reduction buys precision per replication; it does not substitute for replications,
and any ranking reported at that budget is an artifact of it.
\end{quote}

\subsection{Sensitivity, outcome definitions, and scope}

The simulation's parameters were subjected to: global multiplicative jitter at
three amplitudes over 1,002 draws; a tiered design and a randomized-matrix design
at 1,000 draws each; 944 multi-level one-at-a-time perturbations covering all 118
registered values at four distances in each direction; a twenty-trajectory Morris
elementary-effects screen over all 118; a Sobol decomposition on a 1,024-row base
design over the eight Morris membership leaders; four structural variants that
change the model's architecture rather than its numbers; and a resource-list test.
Full specifications are in the book's technical appendix, sections A3 to A9.

\paragraph{Which parameters matter.} The ordering exponent $p$ has the largest
single influence of any parameter on every outcome the multi-level sweep scores.
Against a full-adherence baseline of 0.0431 maintenance, moving $p$ alone by 25 per
cent in each direction swings maintenance from 0.0027 to 0.2373, a range of 5.45
times baseline; the decay rate $\delta_0$ is second at 3.27 and the step-10 speed
third at 2.22. Across the full four-distance ladder the $p$ range is 14.64 times
baseline, $\delta_0$ 12.98 and member heterogeneity 4.69. The twenty-trajectory
Morris screen ranks the same two first on membership $\mu^*$, at 52.44 and 51.21,
followed by churn at 33.48, dropout sharpness at 24.30, exogenous referral at
23.46, heterogeneity at 18.90 and the step-6 and step-12 speeds at 18.54 and 15.87.
Sobol total-order indices on membership, over those eight leaders at plus or minus
25 per cent, agree: $p$ at 0.576 [0.510, 0.643], $\delta_0$ at 0.373 [0.320, 0.428],
dropout sharpness at 0.171, churn at 0.132, heterogeneity at 0.079 and exogenous
referral at 0.058, against a noise floor of 0.043 for membership and 0.073 for
practice. The two step speeds, at 0.029 and 0.038, sit at or below that floor and
are not separated from Monte Carlo error.

At the retired 128-row base sample the first-order column failed three internal
diagnostics and was withheld. The 1,024-row release design repairs the membership
column and not the practice column. On membership no factor now has $S_1 > S_T$,
and the first-order indices sum to 0.693, so the column is admissible: $p$ resolves
at 0.404 [0.288, 0.525] and $\delta_0$ at 0.238 [0.156, 0.327], while the remaining
six have intervals covering zero and are unresolved rather than zero. On practice
the column is still not usable: $\delta_0$ returns $S_1 = 0.421$ against
$S_T = 0.417$, violating the identity $S_T \ge S_1$ that holds for any true
decomposition, and the practice first-order indices sum to 1.074, whereas a sum of
first-order indices cannot exceed one. No practice first-order number is quoted
elsewhere. The sums of $S_T$ are 1.456 for membership and 1.464 for practice, and
the excess over one is the signature of interaction counted once per factor
involved, so interactions are present and are not dominant. The decomposition is
conditional on these eight factors and these ranges; it is not a decomposition of
the model's total variance.

\paragraph{Structural variants.} Four changes to the architecture, not the
numbers, are compared with the base model: a flat per-step gate, admission moved
to Tradition~3, capacity supplied by all members rather than the established
core, and a piecewise-linear rather than hyperbolic member-side capacity response.
The last variant still clips at one; its historical filename is shorthand, not a
literal absence of saturation. Each of the
five architectures is run at 400 paired seeds for five scenarios. Under pure
attraction loss, endpoint existence is 1.000, 0.870, 1.000, 1.000 and 1.000;
endpoint viability is 0.985, 0.658, 0.985, 1.000 and 0.995; and mean membership is
12.38, 6.61, 12.38, 14.73 and 13.04. Under referral loss the corresponding values
are 0.105, 0.120, 0.105, 0.258 and 0.145 for existence; 0.028, 0.055, 0.028, 0.123
and 0.058 for viability; and 0.51, 0.75, 0.51, 2.00 and 1.02 for membership. Thus
the referral-loss ordering holds on all three outcomes in all five tested
architectures. This is a finite structural audit, not a proof over untested
architectures.

\begin{quote}\small
\textbf{What the structural audit says, and no more.} In the corrected model,
referral loss is worse than pure attraction loss on existence, endpoint viability
and mean final membership in the base architecture and all four variants. The
earlier result in which the size ordering reversed under three variants came from
the retired, uncentred capability model and does not reproduce. The parameter
screens remain a separate question and report strict orderings, ties and reversals
at each design level. Specific probabilities and memberships remain conditional
statements about a constructed model.
\end{quote}

\paragraph{What the designs could in principle have found.} Two failures are worth
naming because both were made in this project. A multiplicative perturbation
cannot move a structural zero, so no multiplicative design is evidence about the
two-tier split (Section~4.4). And a $\pm30$ per cent jitter cannot tell you
whether a result depends on the magnitudes at all, only whether it tolerates small
disagreement about them; the structural randomization is the design that answers
the first question, and it is the one the coupling claims fail.

\subsection{Founding composition: an unresolved comparison}

Twenty-five founders with a fixed total practice of 13.75 distributed three ways
(even: all at 0.55; concentrated: five at 1.00 and twenty at 0.4375; split: twelve
at 0.90 and thirteen at 0.2269), 400 paired seeds each. Membership ends at 17.80,
18.09 and 17.45 with half-widths 0.88, 0.99 and 0.93; endpoint viability is 0.985,
0.9875 and 0.985; established practice is 0.2645, 0.2628 and 0.2660. Relative to
even founders, the paired membership differences are 0.29 [-0.90, 1.48] and
-0.35 [-1.47, 0.78]. No equivalence margin was prespecified, so the result is
unresolved rather than evidence of equality.

The unresolved comparison should be discounted heavily before it is read. Founding
practice changes ordering gates, maintenance, resource capacities, dropout and
attraction jointly. Resources are computed from aggregates and no member's state
appears in another member's growth equation except through those aggregates, so the
model has no representation of mentoring, pairing, cliques or sponsorship.
\textbf{It also has no
representation of the mechanism that produced the Carrell--Sacerdote--West
result}, which is people choosing whom to associate with inside a group whose
composition has been arranged. That experiment engineered Air Force Academy
squadrons from measured peer effects, predicted a gain of 0.053 grade points for
the bottom third, and measured a treatment effect of $-0.061$ on exactly the
students it set out to help, because the low-ability cadets re-sorted toward each
other. A non-significant contrast from a design without that mechanism is not
evidence of practical equality.

The design also cannot separate two things: the conditions differ in the variance
of founding practice \emph{and} in the number above the stricter 0.5 experienced
threshold (25, 5 and 12). All 25 exceed the 0.1 established threshold in every
condition. A resolved difference could still have arisen through any of several
state-dependent channels.

\section{The Comparative Case}

A formal correspondence gains little from a single historical case and can lose a
great deal by leaning on one. This section is included because the case bears
directly on the paper's central claim, and because the standard account of it,
which is what a reader is most likely to bring, turns out to be false.

\subsection{The received account, and what the primary source says}

The Washingtonian Total Abstinence Society was founded in Baltimore in April 1840
by six working men, in a scene whose earliest surviving account is Harrison (1860)
and which reaches most modern readers through Maxwell (1950), spread nationally within four years on claims reaching into
the hundreds of thousands, and was effectively finished within a decade. In the
literature descending from AA it is standardly described as a movement that died
of having no rules, and it is used to illustrate why the Traditions matter. That
literature arrives with its moral pre-attached.

The movement's own manual falsifies the account. Grosh (1842), the
\emph{Washingtonian Pocket Companion}, printed at Utica and by then in a second
edition, carries in its first fifteen pages a definition of principles, directions
for organizing and conducting meetings, and a model constitution. Among its
provisions: each society independent and subordinate to none; funds controlled by
its own members; and nothing sectarian or political admitted to lectures,
speeches, singing, or the doings of the society. A footnote records that a
Washingtonian mass convention at Utica passed a declaration of principles and a
model constitution on 22 February 1842, that it was printed in the \emph{Utica
Washingtonian} of 25 February and reprinted in October because of demand, and that
a copy should be procured wherever a society is organized. \textbf{They had
written rules, they had them within two years of founding, and they had a
mechanism for transmitting them.}

What they had, in writing, were analogues of four of the Traditions this analysis
classifies as protective. What they had none of were the seven classified as
enabling. And on anonymity they took the opposite position deliberately and with
an argument: Grosh's directions for a first meeting have joiners rise and call out
their names for the secretary, because ``publicity and freedom are preferable to
private solicitations, whisperings, and secresy in giving the names.''

The comparison the case actually supports is therefore not rules against no rules.
It is one written code against another, and the provisions missing from the first
are the ones the aggregation condition points at.

\subsection{Prominence without anonymity}

The touring-speaker structure of \S3.6 was built for this case. Washingtonian
expansion from 1841 to 1843 proceeded through touring speakers addressing local
societies, which is exactly a rise in the proportion of total attention carried by
a one-directional cross-society channel. By equation~(\ref{eq:speakers}) the
speakers' share of influence is fixed by the ratio of outward to returned
attention and does not fall as the movement grows, so the movement could not have
outgrown the exposure.

The exposure was realized. John B. Gough, the movement's most prominent speaker,
relapsed publicly in September 1845; opponents seized on it, and public confidence
in the movement was impaired. His own account, read at source, contains the words
``I have fallen,'' an acceptance of blame, and a submission to his church's
judgment. It is worth reading rather than inferring from its chapter headings,
which suggest a man rebutting a charge rather than owning a relapse. He had
relapsed once before, twenty-nine months earlier, when his influence weight was
small, and that episode was handled internally and quickly. Same man, same
illness, same candor: what differed was the weight.

\subsection{Maxwell, and a priority problem stated plainly}

Milton Maxwell's 1950 comparison of the two fellowships reaches, without any
formal apparatus, a substantial part of this paper's conclusion. His final section
lists AA's advantages as exclusively alcoholic membership, singleness of purpose,
a definite program of recovery, anonymity, and what he calls hazard-avoiding
traditions; he writes that a comparison with the Washingtonian experience
underscores \emph{the sheer survival value} of the principle of anonymity; and he
reaches that conclusion by way of Gough's relapse and what it cost a movement
whose credibility sat in named men. He also identifies the tradition of keeping
authority in principles rather than in offices and personalities, and connects it
to rotating leadership.

\textbf{What this paper adds to Maxwell is the theorem, and nothing else.} He had
the observation, the mechanism, and the case. He had no formal condition to which
the observation could be referred, and therefore no way to say why anonymity
should have survival value rather than merely that it did. Whether the
vanishing-influence condition is what Maxwell was pointing at is a separate
question, and this paper does not settle it.

\subsection{What the case cannot do}

It is one case, selected because it is the obvious comparison, and the direction of
selection is unfavorable: the Washingtonians are famous among people interested in
AA precisely because the contrast is instructive. Nothing here establishes that the
missing enabling provisions caused the decline; the movement was also absorbed by a
temperance politics it had defined itself against, and the manual's anti-politics
article had nothing to bite on once the movement's identity was itself a political
position. How widely the Utica model constitution was actually adopted is not
recorded in the manual and we have found no source that settles it. The historical
material is offered as an existence proof that the distinction between the two tiers
of rule is visible in a real code, not as evidence about the consequences of
omitting one tier.

\section{Limitations}

\subsection{The largest one, stated first}

\textbf{The mapping in Table~\ref{tab:mapping} is an interpretation of the
Traditions' wording, arrived at by the author, and no computation in this paper
touches it.} Everything the paper claims about AA specifically, as opposed to
about stochastic matrices, passes through it. A reader who holds that Tradition~2
concerns humility rather than weighting, or that anonymity is chiefly protective
of individuals, can accept every number here and reject the paper's thesis. This
is not a caveat on a result; it is the status of the result.

What would settle it is elicitation: give the Traditions' published short text to
readers who do not know the hypothesis, ask them to say what each rule constrains,
and measure agreement with the mapping. That has not been done. A single reader
who knows the rooms did read a draft and identified a substantive error, the
membership-versus-attendance conflation of \S6.1, and no other intervention in the
project changed as much for as few words. The implication for the parts nobody has
checked is uncomfortable and is the reason this limitation is listed first.

\subsection{The remaining limitations}

\begin{enumerate}
\item \textbf{No parameter is estimated.} Section~3's computations demonstrate a
theorem on constructed matrices; Section~6's are simulations from 118 assumed
values, none fitted, because the longitudinal data such a model would need has
never been collected. The paper's empirical content is its predictions, not its
numbers. A model of this kind can show that a set of ideas is consistent and that
a mechanism is available. It cannot show that anything is true.

\item \textbf{The $S$ and $G$ matrices are hand-written judgments}, and all of
Section~4 plus the resource structure of Section~6 inherits from them. Their
tolerance of $\pm30$ per cent perturbation is necessary and not sufficient
support, and the structural randomization they fail (\S4.5) is the design that
speaks to the question. The correct remedy is a second governance matrix elicited
independently from another reader, scored by Cohen's kappa against the first. A
blank elicitation form exists; no second reader has completed it. Until one does,
Section~4 rests on one person's judgment.

\item \textbf{Member heterogeneity was tuned to produce a desired behavior.} Its
standard deviation was selected in an attempt to convert individual cliffs into a
graded group response. After mean-centring and re-estimating the state-dependent
environment, the corrected audit finds typical-member bistability in only 1.75 per
cent of 400 endpoint environments. The intended mechanism is therefore not established
by the released baseline, and Section~6's quantities inherit from the authored spread.

\item \textbf{Specification search.} Several architectural choices (the retention
channel for Tradition~3, core-based carrying capacity, the placement of the
capacity gate) were refined against simulation behavior as well as substantive
reasoning. This risks tailoring an architecture to expected behavior and no
assurance can be given that it has not happened. The four structural variants of
\S6.6 are a partial and inadequate answer.

\item \textbf{Thirty-five parameters cannot affect a fully adherent group at all},
because the governance matrix cancels exactly at full adherence. Any sensitivity
result quoted over the full 118 must be read with that in mind, and a design that
perturbs those 35 and reports no effect has found nothing.

\item \textbf{Conditional figures read as unconditional.} Quality in
Table~\ref{tab:decline} is computed among members of surviving groups. Read
without the surviving fraction beside it, ``quality is maintained'' would describe
a population that has partly ceased to exist. Table~\ref{tab:conditional} exists
to make the difference visible.

\item \textbf{DeGroot averaging is a strong simplification of a group
conscience.} Real members argue, defer selectively, abstain, update out of order,
and sometimes harden rather than converge. Golub and Jackson's result is about
naive averaging specifically. Whether real deliberating groups behave like DeGroot
updaters has not been tested here or, so far as we have found, anywhere.

\item \textbf{Endogeneity of the group input} blocks estimation of Section~5's
technology and is unsolved here.

\item \textbf{The sign of $\rho$ has never been measured.} The contribution is to
make the question answerable, not to answer it.

\item \textbf{No representation of who attends to whom.} In Section~6's model no
member's state enters another member's growth equation except through group
aggregates, so there is no sponsorship, no pairing, no clique, and no
re-sorting. The mechanism that produced the Carrell--Sacerdote--West result is
absent by construction (\S6.7). Building it is a different model, not a different
run, and it is the largest single piece of technical work outstanding.

\item \textbf{Durations were never calibrated.} The horizon, the arrival rate and
the churn floor are three of the 118 registered numeric values. Endpoint orderings
have outcome-specific sensitivity records, and the 10-, 20-, 30-, 50- and
100-year checks continue to move. The trajectories are shapes inside the model,
not forecasts of how long a real group lasts or evidence of steady state.

\item \textbf{Combinations were not tested.} Each decline condition switches one
thing off at full adherence elsewhere. Real decline is unlikely to be so tidy and
there is no reason to expect the costs to add.

\item \textbf{A correspondence is not a cause.} That the Traditions implement a
known aggregation criterion does not establish that this is why they were written,
or why AA has endured.

\item \textbf{Sources not read.} Alexander (1988) on the class and domestic
dimensions of the Washingtonian movement; Blumberg (1980, 1991) on its political
entanglement; the full text of Pagano et al.\ (2004), which is cited at a remove;
White's \emph{Slaying the Dragon}. Maxwell (1950) was read in full, but the copy
available to us is a retyped web reproduction with visible transcription errors,
not a scan of the journal, so every Maxwell citation here has been checked against
a transcription rather than against the journal. \emph{Twelve Steps and Twelve Traditions} (1953) was read in
full on 10 August 2026 and is held as a record with no document; the rest of AA's
own literature has not been obtained.
\end{enumerate}

\section{Falsifiable Predictions}

Where the analysis reported above bears on a prediction, or narrows what it should
say, this is marked.

\begin{enumerate}
\item \textbf{Decision quality improves with group size under flat influence and
plateaus under concentrated influence.} Testable with forecast or vignette tasks
administered to group consciences of varying structure. Unchanged, and untested.

\item \textbf{Rotation breadth, not rotation per se, predicts group durability.}
Testable from service rosters. \emph{Narrowed by the analysis above:} the operational threshold is
roughly a quarter of the membership in the rotation pool, computed as
twenty-six per cent across sizes from 50 to 800 (\S3.5). The effect of falling
short is a real but moderate permanent cost, not a near-clique regime, and a test
should be powered accordingly.

\item \textbf{Groups losing exogenous newcomer inflow decline demographically
while established practice among selected viable remnants can remain high.}
Unlisted meetings, schedule-disadvantaged meetings, and meetings distant from
referral sources provide natural variation. \emph{Narrowed by the corrected
trajectories:} viable-room membership also falls sharply, so the interior is not
warning-free. The testable divergence is between losses in groups and members and
the smaller change in practice among those who remain; it must be measured on all
three series.

\item \textbf{The two supply channels dissociate}: attraction loss produces stable
remnants, referral loss produces delayed dissolution, so proximity to treatment
facilities should predict survival through low-attraction periods. \emph{Partly
answered, and qualified.} Earlier drafts described this as one of exactly two
universal simulation results. The corrected analysis reports strict support,
ties, reversals, and unresolved screens separately; it does not promote a finite
sensitivity suite into a universal claim. Mortality and membership remain
distinct estimands.

\item \textbf{Tradition~3 has distinct resource-governance and inverse-practice
friction paths.} \emph{Bearing of the analysis above:} at 400 paired seeds, combined
loss ends at 6.76 members against 17.80, closes 25.0 per cent of groups, and leaves
54.8 per cent endpoint-viable against 98.5 per cent at baseline. Friction loss alone
costs 2.96 members [1.85, 4.08]; governance loss alone costs 6.03 [5.04, 7.01].
The earlier predominantly-size interpretation does not survive the corrected model.

\item \textbf{If step order does not predict step completion} ($\rho > 0$
throughout), the chain interpretation of the CES reformulation fails. The
proxy-averaging exercise in \S5.1 neither identifies latent practice nor establishes
a required number of instruments.

\item \textbf{the two-tier division should be visible to readers who do not
know the hypothesis.} Give the published short text of the Traditions to
independent readers and ask which of them supply anything a member uses directly
and which only protect other rules. The present split is author-coded, not a tier
discovered by the model. If independent readers do not reproduce it, the resource
layer is encoding the author's prior rather than the text's content.

\item \textbf{the mapping itself is testable by elicitation.} See \S8.1. This
is the prediction whose failure would cost the paper the most.
\end{enumerate}

\section{Conclusion}

This paper's durable contribution is a correspondence: three of AA's Twelve
Traditions, on a reading of their wording that the paper states plainly and does
not verify, jointly implement the vanishing-influence condition under which naive
collective deliberation aggregates information, a criterion formalized by Golub
and Jackson in 2010 and reached by a fellowship of laypeople, from eleven years of
watching groups fail, in 1946. The correspondence yields two corollaries with
immediate empirical content, that rotation must scale with the group and that
obstruction is a matter of scaling rather than severity, and it reframes anonymity
from an ethic of humility to a structural precondition of trustworthy group
decision-making.

Around that center, the paper contributes a method: deriving institutional
couplings through an explicit resource layer rather than asserting them, which
rejects the natural index-pairing conjecture on all twelve of its instances while
recovering, unbidden, the functional division of the Traditions and the primacy of
unity. It also contributes an accounting of what that method cannot support. Five
of the twelve rejections are structurally forced and invisible to the perturbation
design that was originally quoted for them; a threshold test that can see them
finds them failing by comparable margins; and the central coupling claim survives
structural randomization less often than it fails. Those facts are reported here
because a derivation whose robustness is asserted rather than measured is worth
less than one whose limits are known.

The membership model's principal prediction, that referral-starved groups decline
demographically while looking healthy from inside, is specific, mechanistically
grounded, distinguishable from rival accounts, and testable with records AA's
service structure already keeps. In the corrected five-architecture comparison,
referral loss is worse than pure attraction loss on existence, endpoint viability
and mean final membership. The expanded parameter screens bound that result over
their stated ranges; they are not a proof outside them.

Where this line of work should go next is not further modeling. It is
measurement: of the mapping, by elicitation from readers who do not know the
hypothesis; of the governance matrix, by a second independent elicitation; of
rotation breadth against group longevity; of step ordering against step
completion; and of the quiet demographic signature of groups that seekers have
stopped finding.

\section*{References}
\addcontentsline{toc}{section}{References}

\noindent\textit{Status note.} An entry explicitly marked read at source, read in
full, abstract only, cited at a remove, or not read has that current-project
status. For entries without an explicit status, the repository does not document
whether the full work was read; they are background citations and do not upgrade
a load-bearing claim. The complete current status register is
\texttt{research/SOURCES.md}. The separately supplied corpus was worked through
on 9 August 2026; what remains under \texttt{research/staged/} is two journal
articles that are recorded as verified online, are unread because retrieval
returned an access challenge, and are not evidence for this paper. No source
document is committed to the repository: each source is published as a citation,
a rights position, a provenance URL, a SHA-256 and a vocabulary-only
verification index.

\begin{list}{}{\setlength{\leftmargin}{1.5em}\setlength{\itemindent}{-1.5em}\setlength{\itemsep}{2pt}}\small

\item Alexander, R. M. (1988). ```We Are Engaged as a Band of Sisters': Class and Domesticity in the Washingtonian Temperance Movement, 1840-1850.'' \textit{Journal of American History} 75(3): 763-785. \textbf{Not read}; cited for the women's dimension of the movement, which this paper does not develop.

\item Angrist, J. D. (2014). ``The perils of peer effects.'' \textit{Labour Economics} 30: 98-108.

\item Banks, H. T., K. L. Rehm, K. L. Sutton, C. Davis, L. Hail, A. Kuerbis, and J. Morgenstern (2014). ``Dynamic modeling of behavior change.'' \textit{Quarterly of Applied Mathematics} 72: 209-251.

\item Banks, H. T., K. Bekele-Maxwell, R. A. Everett, L. Stephenson, S. Shao, and J. Morgenstern (2017). ``Dynamic modeling of problem drinkers undergoing behavioral treatment.'' \textit{Bulletin of Mathematical Biology} 79: 1254-1273.

\item Ben-Porath, Y. (1967). ``The production of human capital and the life cycle of earnings.'' \textit{Journal of Political Economy} 75(4): 352-365.

\item Blair, H. W. (1888). \textit{The Temperance Movement: or, The Conflict Between Man and Alcohol.} Boston: William E. Smythe. \textbf{Read at source}; public domain.

\item Blumberg, L. U. (1980). ``The Significance of the Alcohol Prohibitionists for the Washingtonian Temperance Societies.'' \textit{Journal of Studies on Alcohol} 41(1): 37-77. \textbf{Not read.}

\item Blumberg, L. U., and W. L. Pittman (1991). \textit{Beware the First Drink! The Washingtonian Temperance Movement and Alcoholics Anonymous.} Seattle: Glenn Abbey Books. \textbf{Not read}; cited as an outstanding book-length comparison, not as evidence for a specific result.

\item Carrell, S. E., B. I. Sacerdote, and J. E. West (2013). ``From Natural Variation to Optimal Policy? The Importance of Endogenous Peer Group Formation.'' \textit{Econometrica} 81(3): 855-882. doi:10.3982/ECTA10168. \textbf{Read at source} from the lead author's university copy. Earlier circulated as NBER Working Paper 16865 and, before that, as \textit{Beware of Economists Bearing Reduced Forms?}. In copyright; full text not redistributed.

\item Crothers, T. D. (1911). \textit{Inebriety: A Clinical Treatise.} Cincinnati: Harvey Publishing. \textbf{Read at source}; public domain.

\item Cunha, F., and J. J. Heckman (2007). ``The technology of skill formation.'' \textit{American Economic Review} 97(2): 31-47.

\item Cunha, F., J. J. Heckman, and S. M. Schennach (2010). ``Estimating the technology of cognitive and noncognitive skill formation.'' \textit{Econometrica} 78(3): 883-931. \textbf{Read at source.}

\item DeGroot, M. H. (1974). ``Reaching a consensus.'' \textit{Journal of the American Statistical Association} 69(345): 118-121. \textbf{Read at source}; the updating model.

\item Eddy, R. (1887). \textit{Alcohol in History.} New York: National Temperance Society. \textbf{Read at source}; public domain.

\item Fehlandt, A. F. (1904). \textit{A Century of Drink Reform in the United States.} Cincinnati: Jennings and Graham. \textbf{Read at source}; public domain.

\item Fatimah, H., M. D. Hunter, and M. A. Bornovalova (2025). ``Modeling the Dynamics of Addiction Relapse Via the Double-Well Potential System.'' \textit{Journal of Psychopathology and Clinical Science} 134(1): 69-80. doi:10.1037/abn0000960. \textbf{Read in full} from the author manuscript; the strongest empirical warrant used here for a two-well relapse landscape, with the limitations stated in Section~5.

\item Galanter, M. (1981). ``The `relief effect': A sociobiological model for neurotic distress and large-group therapy.'' \textit{American Journal of Psychiatry} 138(5): 588-591.

\item Golub, B., and M. O. Jackson (2010). ``Na\"ive Learning in Social Networks and the Wisdom of Crowds.'' \textit{American Economic Journal: Microeconomics} 2(1): 112-149. \textbf{Read at source.} The wisdom criterion, the three obstructions, and the convergence conditions. Every formal claim in Section~3 originates here.

\item Gorman, D. M., J. Mezic, I. Mezic, and P. J. Gruenewald (2006). ``Agent-based modeling of drinking behavior.'' \textit{American Journal of Public Health} 96(11): 2055-2060.

\item Gough, J. B. (1869). \textit{Autobiography and Personal Recollections of John B. Gough.} Springfield, Mass.: Bill, Nichols \& Co. \textbf{Read at source}; public domain. The September 1845 episode in his own words.

\item Grosh, A. B., comp. (1842). \textit{Washingtonian Pocket Companion.} Second edition. Utica, N.Y.: B. S. Merrell. \textbf{Read at source}; Harvard copy digitized by Google, via HathiTrust, \url{https://hdl.handle.net/2027/hvd.32044004487591}. Public domain. The definition of principles, the model constitution's articles, the Utica mass convention of 22 February 1842, and the directions for taking names publicly.

\item Harrison, D., Jr. (1860). \textit{A Voice from the Washingtonian Home.} Boston. \textbf{Read at source}; public domain. The earliest account of the founding scene.

\item Hawkins, W. G., ed. (1862). \textit{Life of John H. W. Hawkins.} Boston: Briggs and Richards, sixth thousand. \textbf{Read at source}; public domain. Note the edition: this is not the Jewett printing usually cited.

\item Holmstr\"om, B. (1982). ``Moral hazard in teams.'' \textit{Bell Journal of Economics} 13(2): 324-340.

\item Hu, Y., and S. M. Schennach (2008). ``Instrumental variable treatment of nonclassical measurement error models.'' \textit{Econometrica} 76(1): 195-216.

\item Hufford, M. R., K. Witkiewitz, A. L. Shields, S. Kodya, and J. C. Caruso (2003). ``Relapse as a nonlinear dynamic system.'' \textit{Journal of Abnormal Psychology} 112(2): 219-227.

\item Humphreys, K., L. A. Kaskutas, and C. Weisner (1998). ``The Alcoholics Anonymous Affiliation Scale.'' \textit{Alcoholism: Clinical and Experimental Research} 22(5): 974-978.

\item Iannaccone, L. R. (1992). ``Sacrifice and stigma: Reducing free-riding in cults, communes, and other collectives.'' \textit{Journal of Political Economy} 100(2): 271-291.

\item Kaskutas, L. A., J. Bond, and K. Humphreys (2002). ``Social networks as mediators of the effect of Alcoholics Anonymous.'' \textit{Addiction} 97(7): 891-900.

\item Kelly, J. F., K. Humphreys, and M. Ferri (2020). ``Alcoholics Anonymous and other 12-step programs for alcohol use disorder.'' \textit{Cochrane Database of Systematic Reviews}, CD012880.

\item Krout, J. A. (1925). \textit{The Origins of Prohibition.} New York: Alfred A. Knopf. \textbf{Read at source}; public domain. Independent corroboration of the founding, the officers, the fee and the dues.

\item Kurtz, E. (1991). \textit{Not-God: A History of Alcoholics Anonymous.} Expanded edition. Center City, Minn.: Hazelden. \textbf{Consulted at source}; in copyright, full text not stored. The Traditions' drafting history. A vocabulary-only verification index is retained at \texttt{research/incorporated/Kurtz\_1991/} in place of the text.

\item Lembke, A. (n.d.). ``Sacrifice, stigma, and free-riding in Alcoholics Anonymous.'' Association for the Study of Religion, Economics and Culture. \textbf{Cited at a remove.}

\item Marsh, J. (1866). \textit{Temperance Recollections.} New York: Charles Scribner. \textbf{Read at source}; public domain.

\item Alcoholics Anonymous World Services (1953). \textit{Twelve Steps and Twelve Traditions.} New York: AAWS. \textbf{Read in full} 10 August 2026, from the per-chapter files AAWS publishes free at aa.org. Source for the elder-statesman structure of Section~\ref{sec:elders} (Tradition 2, printed pp.\ 132-138), the statement that rotating leadership is best and the warning against entrenched power (Tradition 9, pp.\ 174-178), the two functions of anonymity (Tradition 12, pp.\ 187-191), and the unity/singleness-of-purpose distinction (Traditions 1 and 5, pp.\ 129-131 and 151-155). Also the cross-reference count reported in Section~\ref{sec:conjecture}. \textbf{In copyright; no copy is held in the project repository.}

\item Kurtz, E. (c.\ 1984). ``A Talk About the History of Alcoholics Anonymous From the Letters of Bill Wilson.'' Undated recorded talk; transcript restored by historyofrecovery.com. \textbf{Read in full} 10 August 2026. The year is inferred from internal evidence and is given as approximate throughout. Distinct from Kurtz (1979/1991). Wilson's letters are quoted within it from memory and without page citations, so anything attributed to Wilson through this source is at a remove. \textbf{No copy is held.}

\item Maxwell, M. A. (1950). ``The Washingtonian Movement.'' \textit{Quarterly Journal of Studies on Alcohol} 11: 410-452. \textbf{Read in full}, with a caution: the available copy is a retyped web reproduction carrying visible transcription errors, not a scan of the journal. Every citation to Maxwell in this paper has been checked against that transcription rather than against the journal.

\item Ostrom, E. (1990). \textit{Governing the Commons: The Evolution of Institutions for Collective Action.} Cambridge: Cambridge University Press.

\item Pagano, M. E., K. B. Friend, J. S. Tonigan, and R. L. Stout (2004). ``Helping other alcoholics in Alcoholics Anonymous and drinking outcomes.'' \textit{Journal of Studies on Alcohol} 65(6): 766-773. \textbf{Read in full}; NIH author manuscript, PMCID PMC3008319, obtained 10 August 2026. Source for the 40 against 22 per cent abstinence contrast, its independence from meeting attendance, and the authors' 8 per cent helping-rate limitation.

\item Rohr, R. (2011). \textit{Breathing Under Water: Spirituality and the Twelve Steps.} Cincinnati: Franciscan Media. \textbf{Read in full} 10 August 2026. Cited only for its reading of anonymity as confidentiality and for containing no discussion of the Traditions. The copy consulted was an unauthorised posting; the bibliographic record was confirmed independently of it, no copy is held, and nothing is quoted at length. See \texttt{research/incorporated/Rohr\_2011/} for the rights position in full.

\item Riessman, F. (1965). ``The `helper' therapy principle.'' \textit{Social Work} 10(2): 27-32.

\item Rynes, K. N., and J. S. Tonigan (2012). ``Do social networks explain 12-step sponsorship effects?'' \textit{Psychology of Addictive Behaviors} 26(3): 432-439.

\item S\'anchez, F., X. Wang, C. Castillo-Ch\'avez, D. M. Gorman, and P. J. Gruenewald (2007). ``Drinking as an epidemic.'' In K. Witkiewitz and G. A. Marlatt (eds.), \textit{Therapist's Guide to Evidence-Based Relapse Prevention}, 353-368.

\item Schennach, S. M. (2004). ``Estimation of nonlinear models with measurement error.'' \textit{Econometrica} 72(1): 33-75.

\item Sharma, S., and G. P. Samanta (2015). ``Analysis of a drinking epidemic model.'' \textit{International Journal of Dynamics and Control} 3: 288-305.

\item Tonigan, J. S., G. J. Connors, and W. R. Miller (1996). ``The Alcoholics Anonymous Involvement (AAI) scale.'' \textit{Psychology of Addictive Behaviors} 10: 75-80.

\item Greenfield, B. L., and J. S. Tonigan (2013). ``The General Alcoholics Anonymous Tools of Recovery: The Adoption of 12-Step Practices and Beliefs.'' \textit{Psychology of Addictive Behaviors} 27(3): 553-561. \textbf{Read in full}; NIH author manuscript, PMCID PMC3707937, obtained 10 August 2026. Source for the two-factor structure of step-work and the disagreement between instruments on nine of twelve steps.

\item Witkiewitz, K., and G. A. Marlatt (2004). ``Relapse prevention for alcohol and drug problems.'' \textit{American Psychologist} 59(4): 224-235.

\item Witkiewitz, K., and G. A. Marlatt (2007). ``Modeling the complexity of post-treatment drinking.'' \textit{Clinical Psychology Review} 27(6): 724-738.

\end{list}

\medskip
\noindent\textbf{Referenced but not reproduced.} The Twelve Steps and Twelve
Traditions of Alcoholics Anonymous, paraphrased throughout. The text is copyright
Alcoholics Anonymous World Services, Inc.\ and is not reproduced here. \textit{Twelve Steps
and Twelve Traditions} (1953) was read in full on 10 August 2026 from the files
AAWS publishes free at aa.org; no copy is held, and it is paraphrased rather than
quoted at length. AA's other publications (\textit{Alcoholics Anonymous Comes of
Age}, \textit{Pass It On}, the \textit{Grapevine} essays of 1946, and service
pamphlets) have not been obtained.

\medskip
\noindent\textbf{What was not read.} Any work testing whether real deliberating
groups behave like DeGroot updaters. Any literature on peer-group composition in
voluntary mutual-aid settings specifically. Alexander (1988), Blumberg (1980,
1991), Pagano et al.\ (2004) in full, and White's \textit{Slaying the Dragon}.
Jellinek's per-capita consumption estimates, which are quoted through Maxwell and
whose original has not been traced.

\appendix
\section{The Whole Paper in Plain Language}

\textit{This appendix says everything the paper says, without the math. It is
written for a reader with no background in economics or statistics. Where a
natural way of putting something would be wrong, this appendix says so.}

\subsection{What this paper is about}

AA runs on two sets of twelve ideas. The Twelve Steps are for the person. The
Twelve Traditions are for the group. This paper turns both into math and checks
whether the pieces fit together, and along the way it finds that some things
everybody assumes about them are wrong, and some things nobody says about them are
true.

\subsection{The big finding: why ``nobody's in charge'' actually works}

Think about how an AA group makes a decision. There is no boss. People talk it
over until the room agrees. That is called the group conscience.

Now here is a question: when can you trust a decision made that way? In 2010, two
economists proved the answer. A group that decides by talking it out can be
trusted only if no single person's opinion carries a big fixed chunk of the final
answer. Here is why. Everyone's opinion is partly right and partly mistaken. When
lots of opinions get blended evenly, the mistakes point in different directions
and cancel out, and the more people the better. But if one person's opinion always
makes up a third of the result, their mistakes never cancel. Adding more people
does not help. The group just gets more and more sure of an answer that is no more
likely to be right.

Now look at three of the Traditions. Leaders serve the group; they do not run it.
There is no ladder to climb, because service jobs rotate. And everyone is
anonymous: no last names, no job titles, no status. Each of those rules does the
same thing from a different angle: it stops any one person's voice from getting
too heavy. This paper proposes that those mechanisms can reduce stationary
influence concentration, which is the condition in the 2010 proof. The wording
alone does not establish that they do.

\begin{quote}\itshape
The rules written in 1946 are consistent with a condition mathematicians would
formalize 64 years later. That is a present-day mapping, not evidence that AA's
drafters discovered, intended, or implemented the theorem.
\end{quote}

\textbf{One honest warning about that paragraph, and it is the most important
sentence in this appendix.} The proof is real and we have read it. But the step
from ``Tradition 2 says leaders do not govern'' to ``therefore nobody's opinion
carries extra weight in the math'' is our reading of a sentence. Nobody has
checked it. None of the computer work in this paper tests it. Somebody could
reasonably say those Traditions are about humility, not about arithmetic, and this
paper has no answer for them yet. Everything else here is downstream of that one
unverified step.

\subsection{What it looks like when the rule breaks}

Two familiar characters break it. The first is the old-timer whose opinion ends
every discussion, not because there is a vote, but because everyone waits to hear
what they think. The second is the little circle of long-timers who have already
talked it over before the meeting starts.

Here is the alarming part from the math: when this happens, nothing looks wrong.
The group still reaches agreement. The meetings feel fine. There is no argument
and no drama. The group is simply wrong more often, it is wrong confidently, and
getting bigger does not fix it. A broken group conscience feels exactly like a
working one from the inside.

There is a third case, worse than both, that we did not expect. If a small group
of people not only receives most of the attention but gives all of \emph{their}
attention only to each other, then in the long run the rest of the room counts for
literally nothing. The group's judgment becomes exactly the judgment of that
handful, no matter how many other people are in the chairs.

And it does not take a monster to cause the problem. One person holding five per
cent of the room's attention, which is not very much, already nearly doubles how
wrong a large group ends up.

\subsection{Something the Traditions forgot to mention}

The Traditions say to rotate service jobs so nobody becomes a boss. The math
agrees, and adds a catch nobody wrote down. Rotating only works if lots of
different people take turns. Picture a group of 400 where the same twelve people
trade the jobs among themselves forever. Rotation like that leaves a permanent
floor under how much attention the rotating dozen hold, and the floor does not go
down as the group grows. To really work, roughly a quarter of the group needs to
be in the rotation. Any group could check this against its own service list.

\textbf{One thing not to overstate.} It would be easy to say that a group of 400
rotating twelve people behaves almost the same as a group run by a clique. That is
too strong. It sits in between, and it is closer to the healthy case than to the
clique. The honest description is a real but moderate permanent cost.

\textbf{And a bigger lesson underneath it.} The obvious question to ask about a
prominent member is ``how prominent?'' That is the wrong question. It is ``does the prominence
grow when the group grows?'' A clique that keeps to itself in a room of ten but
gets diluted as the room fills is not a problem at all in the long run. A single
person who stays equally prominent no matter how big the room gets is a problem
forever. Severity is not the test. Scaling is the test.

\subsection{The steps don't match the traditions by number}

It is tempting to think Step 1 goes with Tradition 1, Step 2 with Tradition 2, and
so on down the line. It sounds neat. But it does not hold up, and it was never AA
teaching. The Traditions were written about ten years after the Steps, to fix real
problems groups were having.

Instead of matching them by number, we asked a different question for each step:
what does a person actually need from the group to do this step? Then the matching
takes care of itself. Step 5 is telling another human being the worst things you
have ever done. What do you need for that? You need to know it will not leave the
room. That is anonymity, Tradition 12, not Tradition 5. Step 12 is carrying the
message to the alcoholic who still suffers. Tradition 5 says the group's one job is
to carry its message to the alcoholic who still suffers. Same sentence, basically.
Those two go together.

And when we code this for all twelve steps, two things follow from those authored
choices. First, five of the
Traditions (autonomy, no endorsements, paying your own way, no hierarchy, no
outside opinions) turn out to give members nothing directly. Their whole job is to
protect the other seven from being eroded. Those five empty rows were put into the
table by the author; multiplication did not discover them independently. Second,
unity has the largest semantic overlap under the same coding. That is a concise
description of the table, not a fact about AA groups established by data.

\textbf{Now the part that argues against the finding.} We tested how much of this
survives if you disagree with the numbers we made up.

Unity's importance holds up well. Even if you throw away every number we chose and
keep only the pattern of which rule touches which need, unity still comes out on
top about three times in four.

The rest does not hold up nearly as well. Do the same thing to the
``steps-don't-match-traditions'' finding and it only survives about four times in
ten, which is less often than it fails. Worse, five of our twelve pieces of
evidence for it were not really evidence at all: those five Traditions had zeros
in our table by construction, so of course they never matched. We built a
different test that could actually see those five, and they still fail, by about
the same margins as the other seven. So we think the finding is right. But it
rests on our judgment about the numbers, not on a proof, and anyone who thinks we
chose the numbers badly is entitled to reject it. That is worth saying plainly
rather than dressing it up.

\subsection{Can you skip a step? Nobody has ever checked}

Everyone in AA says you cannot skip a step, that you cannot make honest amends
(Step 9) for harms you never wrote down (Step 4). It sounds obviously true. But as
far as we could find, no researcher has ever actually tested whether people work
the steps in order, or whether doing an early step really is what makes a later
one possible.

This paper turns that belief into one measurable number. If the number comes out
one way, the steps are truly a chain: skip a link and everything after it fails.
If it comes out the other way, they are more like a menu, and a strong group can
carry someone past a step they missed. The number might be measured with a
purpose-built longitudinal design. Our limited exercise adds noise only to an
output while supplying the inputs and loadings without error; it shows that
averaging more output proxies can improve resolution. It does not establish how
many instruments a valid latent-variable design needs.

\subsection{What can actually go wrong for a group}

Start with a fact about AA that shapes everything here: the only requirement for
\emph{membership} is a desire to stop drinking, and no group can take that away
from anybody.

\textbf{One thing that is easy to get wrong here.} It is tempting to put this as
``a group cannot close its doors,'' and that is not right. Closed meetings are
ordinary and long-standing, and a closed meeting is about who is in the room, not
about who is a member. Nobody can be thrown out of AA. Somebody can absolutely be
in a room where they are not welcome. In the executable model Tradition~3 affects
both resource governance and inverse-practice-weighted dropout friction; it does
not record tenure and therefore cannot identify literal newcomer retention. The
two paths are now tested separately.

So if a group cannot revoke anyone's membership, what can actually go wrong? Three
different things, and they leave three different marks:

\begin{enumerate}
\item \textbf{The group becomes invisible.} Nobody finds it anymore: it is not
listed, it meets at a bad time, word of mouth dried up. The model says this group
shrinks to a small huddle kept alive by court cards and treatment-center
referrals. It does not die. It just gets small, and the people in it are somewhat
worse off than they would be in a healthy group.

\item \textbf{The referrals dry up.} The treatment center closed; the court program
ended. The group coasts on word of mouth for about ten years, and then, in roughly
two runs out of three, it quietly dissolves.

\item \textbf{The group gets unwelcoming.} Nobody loses their membership, but there
is an in-crowd, and newcomers can feel it. They come once and do not come back.
This group shrinks by about a third and mostly survives. It is smaller than it
should be.
\end{enumerate}

\textbf{How much you can trust that number depends entirely on how many times the
simulation was run.} From ten runs, that middle case looks like four groups in
five dissolving. From four hundred, it is closer to two in three, and a healthy
group settles at about 42 members rather than 45. Ten runs is far too few. A
number produced by a random simulation and quoted to three digits from ten runs is
not a result at all, and it is very easy to print one without noticing.

And here is the strangest thing the model says. In the case that actually kills
groups, the people still in it are doing fine. Their recovery holds up right to the
end, and if anything looks slightly \emph{better} than average, because the ones
who were struggling have already gone. From the inside, that dying group looks like
a healthy group. Nobody in the room feels anything going wrong. If that is true in
the real world, and it is checkable, it means a group cannot rely on how the
meetings feel to know whether it is in trouble. It has to look at the numbers: how
many newcomers came this year, and how many came back.

\textbf{One qualification, because the loose version of this is wrong.} Quality
does not hold up in all three cases. In the invisible case it drops noticeably.
The sharp version of the claim is narrower and more useful: \emph{the failure that
kills groups is the one with no warning signal inside the room}.

\subsection{Does it matter who starts a group?}

We tested three founding-state distributions with the same mean initial practice.
Their paired final-membership differences are imprecise and no equivalence margin
was specified. The correct result is unresolved, not that all three are the same.

You should not believe that result very much, and the reason is a famous
experiment. Researchers at the Air Force Academy measured how much cadets help
each other study, then used those measurements to build squadrons designed to help
the weakest students. Their own model predicted a gain. What actually happened was
a loss of about the same size, on exactly the students they had tried to help. The
reason was that once you put fifteen strong students and fifteen weak ones in a
room with nobody in between, the weak ones stopped mixing with the strong ones and
found each other instead. The measurement was real; building a room out of it
destroyed the thing that had been measured.

Our model has no way for people to choose who they spend time with, so it could
not reproduce the Air Force mechanism. That study is a warning about transport and
endogenous mixing, not a measurement of the upside or downside of arranging an AA
group. No causal claim about AA composition follows from this comparison.

\subsection{The movement that came before, and what it actually did}

AA was not the first fellowship of drunks helping each other stay sober. In 1840,
six men in a Baltimore tavern founded the Washingtonians. It spread across the
country in four years and was essentially finished in ten.

The story usually told, especially in AA circles, is that they died of having no
rules. \textbf{We checked, and that story is false.} Their own handbook from 1842
turned up, and it contains written rules: each society independent and answering to
nobody above it, funds controlled by its own members, and nothing political or
religious allowed into the meetings. They passed a model constitution at a
convention in February 1842, printed it in a newspaper, reprinted it because of
demand, and told every new society to get a copy.

What they did not have was any of the rules on the other side of the ledger, the
ones that give a member something. And on anonymity they took the opposite view on
purpose, with a reason: their handbook tells new joiners to stand up and call out
their names, because publicity is better than ``whisperings and secresy.''

So the real comparison is not rules against no rules. It is one written code
against another, and the ones missing from the first are exactly the ones the math
points at. That is a better comparison than the one we started with, and it is
better because it is true.

Their most famous speaker relapsed in public in 1845, opponents made the most of
it, and the movement's credibility suffered. He had relapsed once before, quietly,
two and a half years earlier, and the institutional response was much smaller.
The same-person contrast is suggestive, but it does not identify influence weight
as the sole cause and the first episode was not literally costless.

\subsection{How much of this should you believe?}

Here is the honest answer, and it is less flattering than the summary above may
have sounded.

The mathematical theorem in A.2 is real and settled. The claim that AA's three
rules are the same thing as the theorem's condition is \emph{our reading of three
sentences}, and nobody has checked it. That is the whole hinge of the paper and it
is the part with the least support.

The finding about rotation follows from the theorem and is solid math.

The step-and-tradition results come from two tables of numbers we made up, and they
do not survive if you throw those numbers away. We say so in the paper rather than
quoting a robustness figure that belongs to a different claim.

Everything from the computer model has never been checked against real people or
real groups. Every number in it was chosen by us, not measured. The sensitivity
exercises are screens over specified ranges, not proofs that claims rest on no
number. Their strict, tied, and reversed outcomes must be reported separately,
and the earlier referral comparison used a mixed Tradition~11 intervention. The
corrected headline compares referral loss with the pure attraction path, while
structural size and endpoint-viability orderings are reported separately.

A model like that can show that ideas fit together. It cannot show they are true.

Treat the group-decision math as settled mathematics, treat its application to AA
as an unverified reading, and treat everything else as sharp questions waiting for
someone to check them. Be suspicious of any idea, ours included, that feels certain
before it has been tested. The claims here that felt most obvious at the outset
are, without exception, the ones that needed the most narrowing.

\begin{quote}\small
One thing this paper must never be used for: judging any individual person's
recovery. It has nothing to say about whether you, or anyone, is ``doing it
right.'' It has never been tested on a single human being. For anything personal,
talk to a sponsor, a doctor, or a counselor.
\end{quote}

\subsection{The one-sentence version}

\begin{quote}\itshape
The rule that nobody in AA is in charge and nobody uses their last name may not be
just humility: read a certain way, it is the exact mathematical condition that
makes a group's decisions trustworthy, written down by people who learned it from
failure sixty-four years before anyone proved it. Whether that reading is right is
the thing still to be checked.
\end{quote}

\section{Reproducibility Note}

Every quantitative claim in this paper is reproducible from the companion
repository without reference to any external dataset, because there is no external
dataset.

\paragraph{Deterministic results.} Sections 3 and 4 are algebra on constructed
matrices. Influence vectors are normalized left dominant eigenvectors; consensus
errors come from the closed form in equation~(\ref{eq:err}); the coupling is a
matrix product; the threshold test is one division per Step. These carry no
sampling error and are exact to the digits printed.

\paragraph{Stochastic results.} Section 6 reports Monte Carlo output. The design
is: Euler integration at $dt = 0.5$ weeks over 1,560 weeks;
\textbf{400 seeds minimum for any published figure}; survival reported with Wilson
intervals; membership and quality with 95 per cent half-widths from the cross-run
standard error; Tradition comparisons under common random numbers at 400 paired
replications. \textbf{Common random numbers reduce variance and do not reduce the
replication count required}: the comparison in Table~\ref{tab:tradition} was
originally run at 30 paired replications, produced a different ranking, and
reversed when rerun at 400.
Numerical checks use 200 paired seeds at $dt=1.0$, 0.5, 0.25 and 0.125 weeks and
200 seeds at 10, 20, 50 and 100 years. These support the finite-horizon numerical
resolution and show continued long-run movement; they do not prove a steady state.

\paragraph{Perturbation designs.} Multiplicative jitter at
$L \in \{0.15, 0.30, 0.50, 0.75\}$, 2,000 draws per level, seed 3. Structural
randomization replacing every non-zero entry with a uniform draw on $[0.05, 1.00]$,
2,000 draws, seed 23. Wilson intervals throughout on $n = 2{,}000$. A
multiplicative design cannot move a structural zero and is not evidence about
sparsity; the structural design is not evidence about where the zeros belong.

\paragraph{Named threats.} Selection effects in the historical case (\S7.4);
conditioning on survival, addressed by printing the surviving fraction beside
every conditional quality figure; specification search on architecture (\S8.2);
the exact cancellation of 35 governance parameters at full adherence; and, largest
of all, the unverified mapping of \S3.2.

\paragraph{Verification.} All figures are asserted against recomputed values in a
notebook executed end to end in a single process, which exits non-zero on any
exception, any failed assertion, or any cell producing no output. A separate
checker verifies consistency \emph{between} documents: that every decimal is
reachable from the notebook, that simulation output states its seeds and an
interval, and that no citation attributes to a source something that source does
not contain. Each of those checks exists because an error of exactly that kind got
as far as a written page before it was caught.

\end{document}
